% 本文件由 LaTeX 排版 Agent 根据有效 translation.json 自动生成。
% author=Apocrypha; work_id=0801; title=十二位圣祖遗训

\chapter{十二位圣祖遗训}

% structure: p0001 source=h1 target=omitted reason=page-title-already-rendered-by-parent

% structure: p0002 source=h2 target=section reason=relative-heading-rank; base=h2

\section{一、—— \Person{勒乌本}遗训：论思念}

1. \Person{勒乌本}遗训抄本，即他在世一百二十五年临终前所嘱托其子者。在\Person{若瑟}死后二年，他患病时，他的儿子们和孙子们聚集来探望他。他对他们说：我的孩子，我快死了，要走上我祖先的道路。当他看见他的弟兄\Person{犹大}、\Person{加得}和\Person{阿协尔}时，对他们说：请扶我起来，我的弟兄们，好让我向我的弟兄和我的孩子们讲述我心中隐藏的事，因为从今以后我的力量衰微了。于是他起来亲吻他们，哭着说：请听，我的弟兄们，请侧耳听你们的父亲\Person{勒乌本}，我所命令你们的事。看，我今天呼求天上的天主为证反对你们，叫你们不可行走在少年无知和淫乱之中，我贪婪地陷入其中，污秽了我父亲\Person{雅各伯}的床榻。因为我告诉你们，主以严重的灾祸击打我的腰胯达七个月之久；若不是我们的父亲\Person{雅各伯}为我向主祈祷，主必定消灭了我。因为我在三十岁时在主眼前行了这恶事，并且七个月之久我病重垂死；此后我在主面前全心悔改达七年之久。我不饮酒和浓酒，肉不入口，也不尝美味，为我的罪哀悼，因为罪大恶极。在\Place{以色列}绝不可如此行。\par

2. 现在请听我，我的孩子们，我在悔改中所见关于七种迷惑之灵的事。从\OriginalTerm{贝里雅耳}{Beliar}那里，有七种灵被赐予对抗人，它们是少年行为的首领；并且在人受造时，有七种灵赐予他，使人在它们之中完成人的一切作为。第一灵是生命的灵，人的整个存有由此被造。第二灵是视觉的灵，贪欲由此而生。第三灵是听觉的灵，教训由此而来。第四灵是嗅觉的灵，由此赋予嗅吸空气和气息的味觉。第五灵是言语的灵，知识由此而来。第六灵是味觉的灵，饮食由此而来；藉着它们产生力量，因为在食物中是力量的根基。第七灵是生育和性交的灵，藉着它，由于贪爱快乐，罪恶也进入；因此它在创造次序中为末，在少年时期为首，因为它充满了无知，这无知引导少年如盲人入坑，又如牲畜坠崖。\par

3. 此外，还有第八灵，即睡眠的灵，人的本性的迷醉和死亡的形像由此被造。在这些灵中混杂着迷惑之灵。第一，淫乱的灵，居在本性和感官中；第二，腹中贪婪的灵；第三，肝和胆中争斗的灵。第四是谄媚和诡诈的灵，使人因过分殷勤而外表好看。第五是傲慢的灵，使人激动而高傲。第六是谎言的灵，在败坏和嫉妒中伪装言辞，并向亲友隐瞒话语。第七是不义的灵，盗窃和偷窃与之同在，使人实现心中的欲望；因为不义藉着诡计与其他灵合作。此外，睡眠的灵，第八灵，与迷惑和幻象相连。如此，每个少年人都灭亡，使他的心智暗昧而不认识真理，既不理解天主的法律，也不服从父亲的劝诫，正如我年轻时也遭遇的那样。\par

现在，孩子们，要爱真理，真理必保守你们。我劝告你们，要听从你们的父亲\Person{勒乌本}。不可注视女人的容貌，也不可私下与有夫之妇交往，更不可参与妇女的事务。因为若不是我看见\Person{彼耳哈}在遮蔽处沐浴，我就不会陷入这大不义。因为我的心注目女人的裸体，不得安眠，直到我行了那可憎之事。因为我们的父亲\Person{雅各伯}与他父亲\Person{依撒格}同去时，我们在\Place{加德尔}，靠近\Place{白冷}的\Place{厄弗辣大}，\Person{彼耳哈}喝醉了酒，赤身睡在她的内室；我进去看见她的裸体，便行了那亵渎之事，并留下她睡着而离去。立刻，天主的使者将我的亵渎之事启示给我们的父亲\Person{雅各伯}，他来了，为我哀悼，并再没有碰她。\par

4. 因此，不可注目妇女的美貌，也不可思想她们的行径；却要以敬畏主的心，以纯朴的心行事，劳碌工作，游走于研读和羊群之间，直到主赐给你祂所愿的妻子，免得你像我一样受苦。直到我父亲去世，我都不敢正视\Person{雅各伯}的面容，或向我的任何弟兄说话，因为我的羞耻；甚至至今我的良心因我的罪而折磨我。我的父亲安慰了我；因为他为我向主祈祷，使主的怒气离开我，正如主所指示给我的。从此以后，我得到保护，再没有犯罪。因此，我的孩子们，要遵守我所命令你们的一切，你们就不会犯罪。因为淫乱是灵魂的毁灭，使它远离天主，并趋近偶像，因为它欺骗心智和悟性，并在时辰未到之前将少年人引入地狱。因为淫乱毁灭了许多人；即使一个人年老或尊贵，它也使他在\OriginalTerm{贝里雅耳}{Beliar}和人子面前成为羞辱和笑柄。因为\Person{若瑟}使自己远离一切妇女，并洁净自己的思想远离一切淫乱，他便在主和人前蒙恩。因为那埃及妇人向他行了许多事，召来术士，给他爱情药，但他灵魂的决意不容忍任何邪恶的欲望。因此我祖先的天主拯救他脱离一切可见与隐密的死亡。因为如果淫乱不能征服心智，\OriginalTerm{贝里雅耳}{Beliar}也不能征服你们。\par

5. 我的孩子们，妇女是有害的；因为她们对男人没有权力或力量，就用外表的方式巧妙地引诱他们到自己身边；她们用力量不能胜过的，就用诡计胜过。况且，天主的使者曾告诉我关于她们的事，教导我说：妇女比男人更容易被邪淫之神制服，她们在心中图谋对付男人；她们藉着装饰先欺骗他们的心智，用眼波注入毒药，然后以行动俘虏他们，因为妇女不能以武力胜过男人。\par

因此，我的孩子们，要逃避邪淫，并命令你们的妻子和女儿不可装饰她们的头部和面孔；因为凡是如此行事诡诈的妇女，都已注定受永罚。因为她们就是这样诱惑了洪水前的\OriginalTerm{守望者}{Watchers}；这些守望者不断地看见她们，就彼此起了欲望，在心中构思行动，并改变成人的形状，在她们与丈夫交合时显现给她们；妇女心中对显现者起了欲望，就生下了巨人，因为守望者显现时身形仿佛直达天上。\par

6. 因此，要谨防邪淫；如果你们愿意心灵纯洁，就要守护你们的感官远离一切妇女。也要命令她们不与男子来往，使她们也保持心灵纯洁。因为经常的会面，即使不行不敬之事，对她们也是不治之症，对我们则是\OriginalTerm{贝利亚}{Beliar}的永久羞辱；因为邪淫本身既无智慧也无虔诚，一切嫉妒都住在其欲望中。因此，你们会嫉妒肋未的子孙，并企图高过他们；但你们不能，因为天主会为他们复仇，你们将死于恶死。因为主将权柄赐给了肋未、犹大、也赐给与我同在的你们，以及丹和若瑟，使我们作统治者。因此我命令你们听从肋未，因为他将知道主的法律，并为全以色列颁布审判和祭献的条例，直到基督——那位主所宣告的大司祭——的时期完成。我指着天上的天主恳求你们，各人与邻人施行真理；并以谦卑的心亲近肋未，好从他口中获得祝福。因为他将祝福以色列；\Emph{特别}是犹大，因为主拣选了他统治万民。我们要敬拜他的后裔，因为他将为我们死于有形和无形的战争，并要在你们中间作永远的君王。\par

7. 勒乌本在给儿子们训示之后便去世了；他们将他放在棺材里，直到从埃及抬上来，葬在赫贝龙的双重洞穴中，那里有他的祖先。\par

% structure: p0012 source=h2 target=section reason=relative-heading-rank; base=h2

\section{II.——西默盎关于嫉妒的遗训。}

1. 以下是西默盎的话，就是他在临终前告诉他的儿子们的，那时他一百二十岁，正是若瑟去世之年。因为他们来探望病中的他，他振作精神坐起来，亲吻他们，对他们说：——\par

2. 我的孩子们，听啊，听你们的父亲西默盎，我心里想什么。我生自我父亲雅各伯，是他第二个儿子；我的母亲肋阿给我起名叫西默盎，因为主俯听了她的祈祷。\BibleRef{创 29:33}我变得极其强壮；我无所畏惧，什么也不怕。因为我的心刚硬，我的意志坚定，我的情感麻木；因为勇力也是至高者赐给人的，在灵魂和身体上。那时我嫉妒若瑟，因为我们父亲爱他；我决意要除灭他，因为欺骗的王子派遣了嫉妒之灵，蒙蔽了我的心智，使我不把他看作兄弟，也不顾惜我父亲雅各伯。但他的天主和他祖先的天主派遣了祂的使者，将他从我手中救出。因为当我前往舍根为羊群取药膏，勒乌本去多堂（那里有我们的必需品和所有物资），我们兄弟犹大将他卖给了依市玛耳人。勒乌本来了就悲伤，因为他本想把他安然无恙地送回父亲那里。但我对犹大生气，因为他让他活着走了，我对他愤恨了五个月；但天主制止了我，使我双手一切工作都受阻，因为我的右手萎缩了七天。我的孩子们，我知道这事是因若瑟临到我身上，我就悔改哭泣；我恳求主恢复我的手，使我免受一切玷污、嫉妒和愚妄。因为我知道我在主和我父亲雅各伯面前对兄弟若瑟图谋了恶事，因为我嫉妒他。\par

3. 现在，孩子们，要提防欺骗和嫉妒之灵。因为嫉妒控制人的整个心智，使人不能吃、不能喝，也不能做任何善事；它总是暗示他除灭他所嫉妒的人；被嫉妒的人却永远兴旺，而嫉妒的人却日渐衰败。有两年之久，我以禁食在敬畏主中折磨我的灵魂，我明白了脱离嫉妒在于敬畏天主。人若逃向主，邪灵就从他身上逃跑，他的心智就变得平安。从此以后，他同情他所嫉妒的人，不谴责那些爱他的人，这样他就停止了嫉妒。\par

4. 我的父亲问起我，因为他看见我忧愁；我便说：\Quote{我的肝疼痛}。因为我比众人更悲哀，因我有份于出卖\Person{若瑟}。当我们下到\Place{埃及}，他把我当作探子捆绑时，我知道自己受罚是公正的，便不忧伤。\Person{若瑟}原是个善人，有天主的圣神在他内：他慈悲怜悯，不向我怀恨；他反而爱我如同爱其余弟兄一样。为此，我的孩子们，你们要提防一切嫉妒和怨恨，以纯朴的心灵和善良的心行走，常常记念你们父亲的兄弟，好使天主也赐给你们恩宠和荣耀，并祝福在你们头上，正如你们在他身上所见。他一生从未因这事责备我们，反而爱我们如自己的灵魂，甚至超过爱自己的儿子；他荣耀我们，并将财富、牲畜和果实白白赐给我们众人。所以，我亲爱的孩子们，你们各人要以善良的心爱自己的弟兄，并从你们中间除去嫉妒的精神，因为这精神使灵魂变得野蛮，毁坏身体；它使人的意图转为愤怒和战争，激起流血，使心智陷入狂乱，不容许谨慎在人身内运作；不仅如此，它还夺去睡眠，引起灵魂的骚动和身体的战栗。因为在睡眠中，有些恶毒的嫉妒迷惑他，啃噬他的灵魂，以邪恶的精神搅扰它，使身体不安，使心智在混乱中从睡眠中醒来；它如同具有邪恶有毒的精神般向人显现。\par

5. 为此，\Person{若瑟}容貌俊美，悦人眼目，因他内里没有存留任何邪恶；因为灵魂的烦扰会显露在脸上。如今，我的孩子们，要在主面前使你们的心善良，在人面前使你们的道路正直，这样你们必在天主和人面前蒙恩。你们要谨慎，不可行淫，因为行淫是万恶之母，使人远离天主，亲近\OriginalTerm{贝里雅耳}{Beliar}。因为我曾在\WorkTitle{哈诺客书}的记载中看到，你们的儿子将与你们一同在淫乱中败坏，并要用刀剑向\Person{肋未}行不义。但他们不能胜过\Person{肋未}，因为他要为主作战，并要征服你们所有的军队；在\Person{肋未}和\Person{犹大}中必有少数人分裂，你们中无人有权柄，正如我父亲\Person{雅各伯}在祝福中所预言的。\par

6. 看，我已将一切事预先告诉你们，为使我对你们灵魂的罪孽无可指摘。如今，如果你们从你们中间除去嫉妒和一切顽固，我的骨头将在\Place{以色列}中如玫瑰绽放，我的血肉在\Person{雅各伯}中如百合花，我的香气将如\Place{黎巴嫩}的香气；如同香柏树，圣者将从我而繁衍，直到永远；他们的枝条将伸展到远方。那时，\Place{客纳罕}的种子必消灭，\Place{阿玛肋克}必无余剩，所有卡帕多西亚人必灭亡，所有赫特人必被全然毁灭。那时，\Person{含}的土地必荒废，所有民族必灭亡。那时，全地从纷扰中得安息，普天下从战争中得安宁。那时，\Person{闪}必受荣耀，因为主天主，\Place{以色列}的大能者，将作为人在地上显现，并由祂拯救\Person{亚当}。那时，所有欺骗的灵都要被践踏，人要管辖恶灵。那时，我要在喜乐中兴起，因祂奇妙的作为而赞美至高者，因为天主取了身体，与人同食，并拯救了人。\par

7. 如今，我的孩子们，要听从\Person{肋未}，并在\Person{犹大}中你们必得救赎；不要对这两个支派骄傲，因为从他们中必为你们兴起天主的救恩。因为主将从\Person{肋未}兴起一位司祭，从\Person{犹大}兴起一位君王，天主而人。祂要拯救所有外邦人和\Place{以色列}家族。为此，我命令你们一切事，为使你们也能命令你们的子孙，叫他们在他们的世代中遵守。\par

8. \Person{西默盎}向儿子们吩咐完毕，就与列祖同眠，享年一百二十岁。他们将他放在一口不朽的木材做的棺材里，为将他的骨骸运往\Place{赫贝龙}。他们在\Place{埃及}人的一场战争中秘密地将骨骸运上去；因为\Place{埃及}人将\Person{若瑟}的骨骸守护在王宫的宝库中；术士们告诉他们，\Person{若瑟}的骨骸离去时，全\Place{埃及}必有黑暗和幽暗，并有一场极大的瘟疫降在\Place{埃及}人身上，以致人连用灯也不能认出自己的弟兄。\par

9. \Person{西默盎}的儿子们按丧礼的律法哀悼他们的父亲，他们在\Place{埃及}直到藉\Person{梅瑟}的手离开\Place{埃及}之日。\par

% structure: p0022 source=h2 target=section reason=relative-heading-rank; base=h2

\section{III.——\Person{肋未}的遗训：论司祭职与傲慢}

1. \Person{肋未}言论的抄本，就是他吩咐他儿子们的一切事，按照他们应行的一切，以及直到审判之日所要临到他们的事。他身体健康时召他们到他跟前，因为已被指示他将要死。当他们聚集时，他对他们说：——\par

2. 我肋未在哈兰受孕，也在那里出生，之后我和父亲来到舍根。那时我大约二十岁，与西默盎一同为我们的妹妹狄纳向哈摩尔报仇。当我们在阿贝耳玛红牧放羊群时，主的一股明悟之灵降临到我身上，我看见所有人都在败坏自己的道路，不义为自己筑起了墙壁，邪恶坐在了塔楼上；我为人类种族悲伤，并向主祈祷，求主拯救我。然后我陷入了沉睡，看见一座高山：这就是阿贝耳玛红的阿斯彼斯山。看，天开了，一位天主的天使对我说：肋未，进来。我从第一层天进入第二层天，看见那里有悬在两层天之间的水。我又看见第三层天，远比那两层天明亮，因为其中有无限的高度。我对天使说：这是为什么？天使对我说：不要为这些事惊奇，因为你将看见另外四层天，比这些更明亮，无可比拟，当你升到那里时；因为你将站在主面前，作祂的仆人，向人们宣告祂的奥秘，并传扬关于那位将要拯救以色列的那一位；\BibleRef{路 24:21}借着你和犹大，主将出现在人间，在他们内拯救所有人类种族；你的生命将是主的产业，祂将是你的田地、葡萄园、果实、金子、银子。\par

3. 那么，请听关于七层天的事。最下层之所以较为阴暗，是因为它接近人所有的不义。第二层有火、雪、冰，为那主所定的日子，在天主正义的审判中预备着；其中有所有报复邪恶的报应之神。第三层有那些为审判之日所预备的军队的众军，要向欺骗和\OriginalTerm{贝里雅耳}{Beliar}的神灵施行报复。直到第四层以上的诸天是圣洁的，因为至高的那一位居住在最上层，在至圣所中，远超一切圣洁。在紧邻的天上有主面前的天使，他们侍奉并在主前为义人的一切无知赎罪；他们向主献上合宜的馨香之气和无血的祭品。在这层天之下，有天使携带答案给主面前的天使。在紧邻的下层天有宝座、宰制者，其中不断向天主献上赞美诗。因此，每当主看我们时，我们全体都震动；是的，诸天、大地和深渊，都在祂威仪面前震动；但人子们不顾这些事，犯罪，激怒至高者。\par

4. 因此，现在要知道：主将审判人子们；因为当岩石崩裂，太阳熄灭，水干涸，火战栗，一切受造物受扰，无形之神灵消融，坟墓因至高者的苦难而被掠夺时，不信的人仍将停留在他们的不义中，因此他们必被刑罚审判。因此至高者听到了你的祈祷，要把你从不义中分别出来，使你成为祂的儿子、仆人和祂面前的仆人。一道明悟的光将照耀在雅各伯中，你将成为以色列所有后裔的太阳。祝福将赐予你和你的所有后裔，直到主藉祂的圣子的慈悲眷顾所有外邦人，直到永远。尽管如此，你的儿子们仍要下手把祂钉在十字架上；因此，你被赐予谋略和明悟，为叫你教导你的儿子们关于祂的事，因为祝福祂的人必蒙祝福，诅咒祂的人必灭亡。\par

5. 天使为我打开天上的门，我看见圣殿和至高者坐在荣耀的宝座上。祂对我说：肋未，我已赐给你司祭职的祝福，直到我来并住留在以色列中间。然后天使带我回到地上，给我一面盾牌和一把剑，说：因狄纳的缘故向舍根报仇，我必与你同在，因为主派遣了我。那时我消灭了哈摩尔的儿子们，正如天上石版上所写的。我对祂说：主啊，我求祢，告诉我祢的名字，好使我在患难之日呼求祢。祂说：我是为以色列种族求情的天使，使祂不至于完全击打他们，因为每个恶灵都攻击它。此后，我如同醒来，赞美了至高者和那为以色列种族并为所有义人求情的天使。\par

6. 当我回到父亲那里时，我发现了一面铜盾；因此那山的名字也叫阿斯彼斯，靠近革巴耳，在阿彼拉的右边；我将这些话记在心里。我与父亲和我的兄弟勒乌本商议，要他命令哈摩尔的儿子们受割礼；因为我嫉妒他们在以色列中所行的可憎之事。我先杀了舍根，西默盎杀了哈摩尔。此后，我们的兄弟们来用刀剑击杀了那城；父亲听说了就发怒，因他们已受了割礼，之后却被处死而忧伤，在他的祝福中他对我们做了不同的事。因为我们违背他的意愿做了这事，我们有罪，那天他病了。但我知道天主的判决是对舍根不利；因为他们想对撒辣行的事，正如他们对我们的妹妹狄纳所行的，但主阻止了他们。他们也迫害我们的祖先亚巴郎，当时他作客旅，并在他羊群增多时抢夺它们；他们羞辱地对待他家中生的仆人耶布拉厄。他们这样对待所有客旅，强夺他们的妻子，并将他们自己驱逐出境。但主的愤怒突然临到他们，直到极致。\par

7. 我对父亲说：\Quote{大人，请勿发怒，因为主必将客纳罕人藉你消灭，并将他们的土地赐给你和你的后裔。从今日起，舍根必被称为愚昧之城；因为人怎样戏弄愚人，我们也照样戏弄了他们，因为他们行愚昧于以色列，玷污了我们的姊妹。然后我们从那里接回我们的姊妹，就出发，来到了贝特耳。}\par

8. 在那里，我又看见一件事，与先前相似，是在我们过了七十天之后。我看见七个身穿白衣的人对我说：\Quote{起来，穿上司铎职的袍子、正义的冠冕、明悟的胸牌、真理的外衣、信德的冠冕、奇迹的冠冕和预言的厄弗得。}他们每人拿着其中一件，给我穿上，说：\Quote{从今以后，你和你后代要永远作主的司铎。}第一个以圣油傅我，并给了我审判的棍杖。第二个以净水洗涤我，并以最神圣的饼和酒喂养我，给我穿上圣洁荣耀的袍子。第三个给我穿上像厄弗得的细麻衣。第四个给我束上紫红色的带子。第五个给我一根丰饶的橄榄枝。第六个在我头上戴上冠冕。第七个在我头上戴上司铎职的冠冕，在我手中放满乳香，使我作为司铎事奉主。他们对我说：\Quote{肋未，你的后裔将分为三支，作为将要来的主荣耀的标记：首先，那忠信者；没有哪一份比他的更大。第二支将在司铎职内。第三支——祂将被称为一个新的名字，因为祂要从犹大兴起作王，并按照外邦人的样式，为所有外邦人建立一个新的司铎职。祂的显现将不可言喻，如同出自我们祖先亚巴郎后裔的一位崇高先知。以色列一切可羡慕的都归于你和你的后裔；一切美观的食物你都要吃；你的后裔将分得主的餐桌，他们中有些将作大司祭、判官和经师，因为圣所将因他们的口被守护。}我醒来后，理解这事与先前相似。我也将这事藏在心里，没有告诉地上任何人。\par

9. 两天后，我和犹大随我们的父亲上到依撒格那里；我祖父按我所见异象的一切话祝福了我，但他不愿与我们同去贝特耳。我们到了贝特耳，我父亲雅各伯在异象中看见关于我，说我应作他们归于主的司铎；他清早起来，藉我将一切什一奉献献给主。我们来到赫贝龙居住，依撒格不断召唤我，使我记起主的法律，正如天主的使者启示给我的。他教导我司铎职的法律、祭献、全燔祭、初熟之果、自愿祭、感恩祭。每日他都教导我，并在主前为我忙碌。他对我说：\Quote{我儿，当谨慎提防邪淫之神；因为这邪淫之神将延续，并藉你的后裔玷污圣物。因此，趁你年轻，为自己娶一个没有瑕疵、未受玷污、不属于培肋舍特人或外邦人族系的妻子。进入圣所前要沐浴；献祭时要洗手；献祭完毕再洗手。从十二种常青的树上，将\Emph{果实}献给主，正如亚巴郎也教导我的；从一切洁净的走兽和飞鸟中向主献祭；从一切头胎和酒中献初果；一切祭品都要用盐调和。}\par

10. 因此，如今要遵守我所吩咐你们的一切，孩子们；因为凡我从我祖先所听见的，我都已告知你们。我远离你们一切的不敬和过犯，就是你们在末世将行反对世界救主的事，不敬虔，欺骗以色列，并从主那里兴起大祸攻击它。你们将不法地对待以色列，以致耶路撒冷不能容忍你们的邪恶；圣殿的帐幔将被撕裂，不再遮盖你们的羞愧。你们要像俘虏一样分散在外邦人中，成为羞辱、咒诅和被人践踏。因为主所拣选的殿宇应称为耶路撒冷，正如义人\Person{哈诺客}的\WorkTitle{书}中所记载的。\par

11. 因此，我娶妻时二十八岁，她的名字是默耳加。她怀孕生了一个儿子，给他起名叫革尔雄，因为我们在本地是寄居的：革尔雄解释为\Quote{寄居}。我看见关于他，他将不在首位。刻哈特在我三十五岁时出生，朝向东方。我在异象中看见他站在全会众的高处。因此我给他起名叫刻哈特，意思是尊威和训导的开端。第三，她在我的四十岁生给我默辣黎；因为他母亲生他困难，就叫他默辣黎，意思是我的苦楚，因为他也死了。约革贝得在我六十四岁时生于埃及，因为那时我在弟兄中出名。\par

12. 革尔雄娶妻，给他生了娄默尼和史米。刻哈特的儿子：阿默兰、依撒阿尔、赫贝隆、乌齐耳。默辣黎的儿子：摩里和贺慕息。在我九十四岁时，阿默兰娶了我的女儿约革贝得为妻，因为他们——他和我女儿——是同日生的。我八岁进入客纳罕地，十八岁杀死舍根，十九岁作了司铎，二十八岁娶妻，四十岁下到埃及。看，你们是我的孩子，我的孩子甚至\Emph{第三代}。在我一百一十八岁时，若瑟死了。\par

13. 现在，我的孩子们，我命令你们：要用全心敬畏我们的主，并按照祂的一切法律，在纯朴中行走。你们也要教导你们的儿女学习，使他们在整个生命中能有理解，不断阅读天主的法律；因为凡认识天主的法律的人，必受尊荣，而无论往何处去，都不会成为外人。是的，他要获得许多朋友，胜过他的祖先；许多人将渴望服事他，并从他的口聆听法律。我的孩子们，要在世上行事正义，好使你们在天上找到\Emph{宝藏}，并在你们的灵魂中播下善事，好使你们在生命中寻获它们。因为如果你们播下恶事，就要收获一切烦恼和痛苦。要殷勤地在敬畏天主中获得智慧；因为纵然人口被掳、城市被毁、土地、金银和一切财产都灭亡，智慧人的智慧却无人能夺去，除非不虔敬的瞎眼和罪的瘫痪：因为即使在仇敌中，他仍显荣耀；在异乡成为家园；在敌人中成为朋友。若有人教导这些事并实行它们，他就要与君王一同被高举，如同我们的兄弟若瑟一样。\par

14. 现在，我的孩子们，我从哈诺客的著作中得知，最后你们将不虔敬地行事，以一切恶意举手攻击主；你们的兄弟将因你们而蒙羞，成为所有外邦人的笑柄。因为我们的祖先以色列将因大司祭的不虔敬而保持纯洁，这些大司祭将举手攻击世界的救主。天是纯净的，高于大地；你们是天的光，如同太阳和月亮。如果你们因不虔敬而黑暗，所有外邦人将做什么？这样，你们就给我们的民族带来诅咒，世界的光就是为这民族而来的，这光在你们中间赐下，为光照每一个人。你们竟想杀害祂，教导违反天主规诫的诫命。你们要抢夺上主的祭品，窃取祂的一份；在向上主献祭之前，你们要取最好的部分，傲慢地与娼妓一同吃。在纵欲中，你们要教导上主的诫命；你们要玷污有丈夫的妇女，亵渎耶路撒冷的贞女；与娼妓和淫妇结合。你们要娶外邦人的女儿为妻，用不合法的洁净来洁净她们；你们的结合将如索多玛和哈摩辣一样不虔敬。你们要因司祭职而骄傲自大，对人自高。不仅如此，还要自高于天主的命令，嘲笑圣物，傲慢地戏弄。\par

15. 因此，主所拣选的圣殿将在不洁中荒凉，你们要在万国中为俘虏，成为他们中的可憎之物；你们要从天主的公义审判中承受凌辱和永远的羞耻；凡看见你们的，都要从你们面前逃跑。若不是为我们的祖先亚巴郎、依撒格和雅各伯，我的子孙中连一个也不会留在地上。\par

16. 现在，我在\WorkTitle{哈诺客书}中得知，你们将迷失七十个星期，亵渎司祭职，玷污祭品，败坏法律，轻视先知的话。在悖逆中，你们要逼迫义人，仇恨虔诚者；厌恶忠信者的言语，称那藉至高者的能力更新法律的人为欺骗者；最后，如你们所想，你们要杀害祂，不明白祂的复活，邪恶地将无辜的血归在自己头上。因祂的缘故，你们的圣所要荒凉，污秽到底，你们没有洁净之处；但在外邦人中，你们要成为诅咒和分散，直到祂再次垂顾你们，并藉着信德和水，在怜悯中接纳你们到自己那里。\par

17. 你们既已听到关于七十个星期的事，也要听关于司祭职的事：因为在每个禧年，都有一个司祭职。在第一个禧年，第一位被傅油进入司祭职的将是伟大的，要如同对父亲一样对天主说话；他的司祭职充满对主的敬畏，在他喜乐的日子，他要起来为世界的救恩。在第二个禧年，被傅油的那位将在所爱之人的忧伤中成孕；他的司祭职要受尊荣，在众人中被荣耀。第三位司祭要在忧伤中被紧紧抓住；第四位要痛苦，因为极大的不义加在他身上，全以色列各人憎恨邻人。第五位要在黑暗中被抓住，第六位和第七位也一样。在第七位中，有如此污秽，以致我不能在主持与人面前表达，因为行这些事的人将知道它们。因此，他们要被掳掠成为掠物，他们的土地和财产要被毁灭。在第五个星期，他们要回到他们荒废的国土，重建上主的殿。在第七个星期，司祭们要来到，他们是拜偶像的、好争的、贪财的、骄傲的、无法无天的、淫荡的、虐待儿童和牲畜的。\par

18. 他们受到主的惩罚之后，主将兴起一位新的司祭，归入司祭职；一切主的话都要启示给祂；祂将在时日的圆满中，在地上施行真理的判断。祂的星要在天上出现，如同一位君王，在日光普照中散发知识之光；祂将在世上被显扬，直到祂被举起。祂要像太阳照耀大地，驱散天下的一切黑暗，普世将有和平。诸天要在祂的日子欢乐，大地要喜乐，云彩要欢腾；主的知识要像海水一样倾注在地上；主荣耀面前的众天使也要在祂内喜乐。诸天要敞开，从荣耀的圣殿中，圣化要藉着圣父的声音临到祂，如同从亚巴郎（依撒格的父亲）。至高者的荣耀要述说在祂身上，智慧和圣化的圣神要在水中安息在祂上。祂要将主的威严永远真实地赐给祂的儿子们；世世代代，直到永远，没有谁能接续祂。在祂的司祭职中，一切罪恶都要终结，不法者要停止作恶，义人要在祂内安息。祂要开启乐园的门，挪开那把守的剑，不让亚当接近；祂要将生命树赐给祂的圣徒吃，圣神要临在他们身上。贝里雅尔要被祂捆绑，祂要赐给祂的子女践踏恶灵的权柄。主要因祂的子女而喜乐，主将永远悦纳祂的爱者。那时亚巴郎、依撒格和雅各伯要喜乐，我也要喜乐，所有圣徒都要穿上喜乐。\par

19. 如今，我的孩子们，你们都已听见；所以，你们为自己选择吧，或是黑暗，或是光明，或是主的法律，或是贝里雅尔的作为。我们回答父亲说：我们要在主面前遵行祂的法律。父亲说：主是见证，祂的众天使是见证，我是见证，你们也是见证，关于你们口中的话。我们说：我们是见证。这样，肋未便结束了向儿子们的训诫；他伸直了脚，被归到他的列祖那里，享年一百三十七岁。他们把他放在棺材里，后来将他葬在赫贝龙，在亚巴郎、依撒格和雅各伯旁边。\par

% structure: p0042 source=h2 target=section reason=relative-heading-rank; base=h2

\section{IV.——犹大关于刚毅、贪财和邪淫的遗训}

1. 犹大在他死前对他儿子们所说的话的抄本。他们聚集到一起，来到他面前，他对他们说：我是我父亲所生的第四个儿子，我母亲给我起名叫犹大，说：我感谢主，因为祂又赐给我第四个儿子。我自幼敏捷活泼，凡事顺从父亲。我孝敬我的母亲和我母亲的妹妹。当我成年时，我的父亲雅各伯为我祈祷说：你将成为君王，在一切事上顺利。\par

2. 主在田间和家中一切工作上恩待了我。我见自己能追上母鹿，便捉住它，为我父亲预备肉食。我猎取瞪羚，也跑赢平原上的一切。我追上野马，捉住并驯服了它；我杀死一头狮子，从它口中夺回一只山羊羔。我抓住熊的掌，将它滚下悬崖；若有什么野兽向我扑来，我就像撕狗一样将它撕开。我遭遇野猪，在追逐中追上它，将它撕裂。在赫贝龙，一只豹子扑向狗，我抓住它的尾巴，将它甩出去，它摔碎在迦萨的边界。一头野牛在田间吃草，我抓住它的角，旋转它使之晕眩，然后扔出去，杀了它。\par

3. 当两位迦南人的王率领大军来攻击我们的羊群时，我独自冲向苏尔王并抓住了他；我击打他的腿，将他拖倒，这样杀了他。另一位王塔弗埃，我杀他时他正骑在马上，于是我驱散了所有百姓。阿科尔王，巨人身材，骑在马上前后投掷标枪，我杀了他；因为我投掷了一块六十斤重的石头，击中他的马，杀死了他。我与阿科尔搏斗了两个小时，杀了他；我将他的盾牌劈成两半，砍断了他的脚。当我剥他的胸甲时，他的八个同伴开始与我搏斗。于是我将衣服缠在手上，用石头打他们，杀了他们四人，其余逃跑了。我的父亲雅各伯杀了贝利撒，他是所有王中的王，巨人般的力量，身高十二肘；恐惧降临在他们身上，他们便停止与我们作战。因此，当我在弟兄们中间时，我父亲在战争中无忧无虑。因为他在异象中关于我见到一位大能的天使时常跟随我，使我不被击败。\par

4. 在南方，我们遭遇了比舍根更大的战争；我与我的弟兄们列阵，追击一千人，杀了他们二百人，以及四位王。我登城攻击他们，又杀了另外两位王；这样我们解放了赫贝龙，夺取了诸王的所有俘虏。\par

5. 次日我们前往阿勒塔，一座坚固、有城墙、无法接近的城，以死亡威胁我们。因此我和加得从城的东边靠近，勒乌本和肋未从西边和南边。那些在城墙上的人以为我们单独前来，便向我们冲下；于是我们的弟兄们悄悄地从两边用梯子爬上城墙，进入城中，那些人并不知道。我们用刀剑夺取了城；那些逃到塔里的人——我们放火烧了塔，连塔带人都拿下了。当我们离开时，塔弗的人攻击我们的俘虏，我们便带着我们的儿子们夺取了他们，与他们交战直到塔弗；我们杀了他们，烧了他们的城，夺取了其中所有的一切。\par

6. 当我在\Place{库泽巴}水边时，\Place{约贝耳}人前来攻击我们，要与我们作战，我们就与他们交战；我们杀了他们来自\Place{塞隆}的盟友，没有给他们任何逃脱或前来攻击我们的机会。第五天，\Place{玛基尔}的人前来抢夺我们的俘虏；我们攻击他们，在激烈的战斗中战胜了他们：因为他们是一支强大的队伍，我们趁他们还未登上山坡就杀了他们。当我们来到他们的城时，他们的女人从城所在的山顶向我们滚下石头。我与\Person{西默盎}藏在城后，占据了高处，彻底摧毁了整座城。\par

7. 第二天，有人告诉我们那两个王的城要带领大军来攻打我们。因此我与\Person{丹}假扮成\Place{阿摩黎}人，作为盟友进入他们的城。深夜，我们的兄弟来了，我们为他们打开城门；我们消灭了所有男子及其财物，夺取了他们的一切，并拆毁了他们的三道城墙。我们逼近\Place{塔姆纳}，那里是敌对诸王的一切避难所。那时我受了伤，便发怒，冲向他们到山顶；他们用石头和标枪投掷我；若非我的兄弟\Person{丹}帮助我，他们本可杀死我。我们于是愤怒地攻击他们，他们都逃跑了；他们从另一条路经过，恳求我的父亲，他就与他们讲和，我们未伤害他们，而是与他们订立了休战协议，并将所有俘虏归还给他们。我建造了\Place{塔姆纳}，我的父亲建造了\Place{兰巴耳}。这场战争发生时我二十岁，\Place{客纳罕}人都惧怕我和我的兄弟们。\par

8. 此外，我有很多牲畜，我的牧人首领是\Person{依兰}\BibleRef{创 38:1}，那个\Place{阿杜兰}人。我去找他时，见到了\Place{阿杜兰}王\Person{巴尔桑}，他为我们设宴；他款待了我，并将他的女儿\Person{巴特叔亚}嫁给我为妻。她为我生了\Person{厄尔}、\Person{敖难}和\Person{舍拉}；其中两个被主击打，无子而死；因为\Person{舍拉}活着，他的子孙就是你们。\par

9. 我们从\Place{美索不达米亚}（从\Person{拉班}那里）来之后，我们的父亲和我们与他的兄弟\Person{厄撒乌}及他的儿子们一起平安地住了十八年。十八年过后，在我四十岁时，我父亲的兄弟\Person{厄撒乌}带领众多强盛的人前来攻击我们；他倒在\Person{雅各伯}的弓下，被抬上\Place{色依尔}山而死；他正是在上了\Place{依辣玛}以上之地时被杀的。我们追赶\Person{厄撒乌}的儿子们。他们有一座铁墙铜门的城；我们无法进入，就扎营围困他们。二十天后他们仍不开门，我当着众人架起梯子，用盾牌护头爬上去，遭受三\Quote{塔冷通}重的石头攻击；我爬上去，杀了他们中四个勇士。第二天，\Person{勒乌本}和\Person{加得}进去杀了另外六十人。于是他们向我们求和；因知晓我们父亲的意图，我们接受他们为进贡者。他们给了我们二百\Quote{苛尔}小麦、五百\Quote{巴特}油、一千五百\Quote{默则}酒，直到我们下到\Place{埃及}。\par

10. 此后，我的儿子\Person{厄尔}娶了来自\Place{美索不达米亚}的\Person{塔玛尔}为妻，她是\Person{阿兰}的女儿。\Person{厄尔}是邪恶的，他怀疑\Person{塔玛尔}，因为她不是\Place{客纳罕}地的人。第三天夜里，一位主的天使击打了他，他未曾认识她，这是由于他母亲的恶谋，因为他不想从她得孩子。在婚筵期间，我将\Person{敖难}许配给她；他也同样在邪恶中未曾认识她，虽然他与她同住了一年。当我威胁他时，他与她同寝……照他母亲的命令，他也死在自己的邪恶中。我想把\Person{舍拉}也给她，但我的妻子\Person{巴特叔亚}不允许；因为她对\Person{塔玛尔}怀有怨恨，因为塔玛尔不是\Place{客纳罕}女子，不像她自己这样。\par

11. 我知道\Place{客纳罕}种族是邪恶的，但青春的意念蒙蔽了我的心。当我见她倒酒时，在酒醉中我被迷惑，在她面前跌倒。我不在时，她去为\Person{舍拉}从\Place{客纳罕}地娶了一个妻子。我知道她所做的，就在灵魂的痛苦中诅咒了她，她也死在她儿子的邪恶中。\par

12. 此后，\Person{塔玛尔}守寡，两年后她听说我要上去剪羊毛；于是她打扮成新娘模样，坐在城对面的城门口。因为\Place{阿摩黎}人有一条法律：将要结婚的女子要在城门口行淫七日。我那时在\Place{科则布}水边喝醉了，因酒未认出她；她的美貌通过她的装扮迷惑了我。我转向她，说：我要与你同寝。她对我说：你给什么呢？我将我的杖、腰带和王冠给了她；我与她同寝，她就怀了孕。我当时不知道她做了什么，想要杀她；但她暗中送还我的抵押物，使我羞愧。当我叫她时，我也听见了我在醉中与她同寝时说的隐秘话；我不能杀她，因为这事出于主。因为我心里说：恐怕她是诡诈行事，从别的女人那里得了抵押物；但我直到死未再亲近她，因为我在全\Place{以色列}行了这可憎之事。此外，城里的人说城中没有新娘，因为她从别处来，在城门口坐了一会儿，她以为无人知道我曾进去找她。此后，因饥荒我们下到\Place{埃及}\Person{若瑟}那里。那时我四十六岁，在那里活了七十三年。\par

13. 如今，我的孩子们，无论我在何事上命令你们，都要听从你们的父亲，遵守我的一切言语，以履行主的诫命，并服从主天主的命令。不要随从你们的私欲，也不要随从你们心中的高傲意念；也不要以青年力壮的事工自夸，因为这也看在天主眼里是邪恶的。因为连我也曾自夸在战争中从未被任何美貌女子的面容所迷惑，并曾责备我的兄弟勒乌本有关我父亲妻子彼耳哈的事，嫉妒和淫乱之灵在我内里排列，直到我跌倒于客纳罕人巴特叔亚和已许配给我儿子们的塔玛尔面前。我对岳父说：\InnerQuote{我要与我父亲商议，然后才娶你的女儿。}他为女儿向我展示了无尽的黄金，因为他是个国王。他用黄金和珍珠装饰她，并在宴会上让她以女子的美貌为我们斟酒。酒使我眼睛迷乱，快乐蒙蔽了我的心；我爱上了她，跌倒了，违反了主的诫命和我祖先的诫命，娶她为妻。主按我内心的意念报应了我，使我在她的儿女身上毫无喜乐。\par

14. 如今，我的孩子们，不要醉酒；因为酒使心意远离真理，在心中点燃情欲之火，使眼睛误入歧途。因为淫乱之灵有酒作为仆役，为心意提供快乐；这两者夺去人的能力。若人饮酒至醉，他便以污秽的念头扰乱心意，趋向淫乱，并刺激身体寻求肉体的结合；若欲望的原因出现，他便行那罪，并不羞愧。我的孩子们，酒就是如此；因为醉酒的人不尊敬任何人。看，连我也被它迷惑，以致我在城中的众人面前不以为耻，因为我当众转向塔玛尔，犯了大罪，揭示了我儿子们的羞耻。我饮酒后，不再敬畏天主的诫命，娶了一个客纳罕女子为妻。因此，我的孩子们，饮酒的人需要谨慎；饮酒的谨慎在于：人应在保持端庄的限度内饮酒；若超过此限，欺骗之灵便会攻击他的心意，成就其意愿；它使醉酒者口出污言，违命而不觉羞愧，甚至以其耻辱为荣，自以为是。\par

15. 凡行淫乱、揭示自己裸体者，已成为淫乱的奴隶，无法脱离其权势，正如我也曾被揭示。我交出了我的杖，即我支派的支柱；我的腰带，即我的能力；我的冠冕，即我王国的荣耀。此后我为此悔改，直到老年不再饮酒或吃肉，也不以任何喜乐为乐。天主的天使向我显示，女人永远统治国王和乞丐；她们从国王夺走荣耀，从勇士夺走力量，从乞丐夺走那支撑他贫穷的一点点。\par

16. 因此，我的孩子们，在酒上要节制；因为其中有四个邪灵：情欲、愤怒、喧闹、不义之财。若你们以欢乐、谦逊和敬畏天主的心情饮酒，你们将得以存活。若你们不以谦逊饮酒，敬畏天主离开你们，那么醉酒便来临，无耻潜入。但\Emph{即使}你们完全不饮酒，也要谨慎，免得在暴怒、斗殴、诽谤和违反天主诫命的言语中犯罪；这样你们将提前灭亡。此外，酒向外人揭示天主和人的奥秘，正如我也曾向客纳罕人巴特叔亚揭示了天主的诫命和我父亲雅各伯的奥秘，而天主禁止我向她宣告这些。酒也是战争和混乱的起因。\par

17. 因此，我的孩子们，我嘱咐你们：不要贪爱金钱，也不要注视女子的美貌；因为为了金钱和美色，我被诱惑而走向客纳罕人巴特叔亚。因我知道，为了这两件事，你们这些我的后裔将陷入邪恶；它们甚至败坏我子孙中的智者，并使犹大王国衰微，这王国是主因我服从父亲而赐给我的。因为我从未违背我父亲雅各伯的一句话，凡他所吩咐的，我都行了。我祖父亚巴郎祝福我，使我成为以色列的王，依撒格也同样祝福我。我知道，王国将从我而建立。\par

18. 因为我也在义人哈诺客的书上读到，你们在末后的日子将行何等邪恶。因此，我的孩子们，要提防淫乱和贪爱金钱；要听从你们的父亲犹大，因为这些事使你们远离天主的法律，蒙蔽灵魂的理智，教导傲慢，不容人对邻人有怜悯；它们夺去灵魂的一切美善，将其捆绑于困苦和患难中，夺去睡眠，吞噬肉体，阻碍天主的祭献；他不记得祝福，不听先知说话，厌烦虔诚的话语。因为一个侍奉两种与天主诫命相悖的情欲的人，不能服从天主，因为它们蒙蔽了他的灵魂，他在白日行走如同黑夜。\par

19. 我的孩子们，贪爱金钱导向偶像；因为人被金钱迷惑时，便提及那些不是神的神明，使拥有它的人陷入疯狂。为了金钱，我失去了我的儿女；若不是我肉身的悔改、灵魂的谦卑和我父亲雅各伯的祈祷，我早已无子而终。但我祖先的天主，慈悲而怜悯，宽恕了我，因为我是在无知中行的。因为欺骗之君蒙蔽了我，我像一个人、像血肉之躯，在罪恶中败坏无知；当我想自己是不可战胜时，却认识了自己的软弱。\par

20. 所以，我的孩子们，你们要知道有两种灵侍候在人旁边——真理之灵和谬误之灵；在两者中间的是明悟之灵，它属于心思，能随意转向任何一方。真理的工作和谬误的工作都写在人的胸前，主知道其中每一个。人的行为没有一时能向祂隐藏，因为在祂面前，人的行为已经写在肋骨的胸上。真理之灵作证一切，控告一切；犯罪的人被自己的心焚烧，不能抬头面向审判者。\par

21. 现在，我的孩子们，要爱\Person{肋未}，好使你们得以存留，不要高抬自己反对他，免得你们被完全消灭。因为主将王国赐给了我，将司祭职赐给了他，祂使王国隶属于司祭职。祂将地上的事物赐给了我，将天上的事物赐给了他。正如天高于地，天主的司祭职也高于地上的王国。因为主拣选了他，超越你们，使他亲近祂，吃祂的餐桌和初果，就是以色列子民中精选的东西，而你们对他们将如同大海。因为，如同在海中，义人和不义的人被抛来抛去，有的被掳去，有的反而致富，同样，在你们内将有人类的每一种族，有的遭险被掳，有的却因掠夺而致富。因为统治的人将如同巨大的海怪，吞吃人如同鱼一般：他们奴役自由的儿女，掠夺房屋、田地、羊群、金钱，用许多人的肉体违法地喂饱乌鸦和鹤；他们将在邪恶中继续前进，在贪婪中不断升级。将有假先知如同暴风雨，他们迫害所有义人。\par

22. 主要使你们彼此分裂，在以色列中常有争战；我的王国将在外邦人中被终结，直到以色列的救恩来到，直到正义天主的显现，使\Person{雅各伯}和所有外邦人安息在和平中。祂要永远护卫我王国的力量；因为主曾向我起誓，说王国必不从我断绝，也不从我的后裔断绝，直到永远，直到永远。\par

23. 我的孩子们，我现在极为忧伤，因为你们的淫乱、行巫术和拜偶像，你们将用来对抗王国，随从那些有交鬼之灵的人；你们将使你们的女儿成为歌女和妓女，用于占卜和谬误的魔鬼，你们将混杂在外邦人的污秽中；因此，主将把饥荒、瘟疫、死亡和刀剑降在你们身上，有报仇的围困，有狗撕裂仇敌，有朋友的辱骂，毁灭和眼目昏花，儿女被屠杀，妻子被掳去，财产被掠夺，天主的殿被焚烧，你们的土地荒凉，你们自己被掳到外邦人中，他们中有些人为他们的妻子阉割你们；每当你们怀着谦卑的心回转归向主，悔改并遵行天主的一切诫命，那时，主将以慈悲和爱眷顾你们，从仇敌的奴役中领出你们。\par

24. 这些事以后，必有一星从\Person{雅各伯}中和平地兴起，有一个人从我的后裔中兴起，如同正义的太阳，以温和和正义与世人同行，在祂内找不着罪恶。天要在他上面敞开，倾注来自圣父的圣神的祝福；祂要将恩宠之灵倾注在你们身上，你们要成为祂真实中的儿子，你们要遵行祂的诫命，从前的和最后的。这就是至高天主的嫩枝，这就是一切血肉之生命之泉。那时，我王国的权杖要发出光芒，从你们的根中要生出一根茎；在其中要生出一根正义之杖给外邦人，审判并拯救所有呼求主的人。\par

25. 这些事以后，\Person{亚巴郎}、\Person{依撒格}和\Person{雅各伯}要复活得生命，我和我的弟兄们将成为首领，就是你们在以色列中的权杖：\Person{肋未}第一，我第二，\Person{若瑟}第三，\Person{本雅明}第四，\Person{西默盎}第五，\Person{依撒加尔}第六，如此按着次序。主祝福了\Person{肋未}；临在天使祝福了我；荣耀的诸能祝福了\Person{西默盎}；天祝福了\Person{勒乌本}；地祝福了\Person{依撒加尔}；海祝福了\Person{则步隆}；山祝福了\Person{若瑟}；帐幕祝福了\Person{本雅明}；天上的光体祝福了\Person{丹}；地的肥油祝福了\Person{纳斐塔里}；太阳祝福了\Person{加得}；橄榄树祝福了\Person{阿协尔}：这样，必有一个主的人民，一种语言；\OriginalTerm{贝里雅耳}{Beliar}的欺骗之灵将不再存在，因为他要被抛入永远的火中。那些在忧愁中死去的人要在喜乐中复活，那些为天主的名而贫穷生活的人要变成富有，那些匮乏的人要得饱足，那些软弱的人要变得强壮，那些为天主的名而被处死的人要在生命中醒来。\Person{雅各伯}的鹿要在喜乐中奔跑，以色列的鹰要在欢欣中飞翔；但不敬神的人要哀叹，罪人全都哭泣，所有的人民要永远赞美主。\par

26. 因此，我的孩子们，要遵守主的一切法律，因为对所有正确遵行祂道路的人，是有希望的。他对他们说：我今天在你们眼前去世，享年一百一十九岁。不要让任何人用昂贵的衣服埋葬我，也不要剖开我的肚腹，因为那是君王们所行的；你们要带我上到\Place{赫贝龙}去。\Person{犹大}说了这些话以后，就安眠了；他的儿子们按他吩咐的一切做了，将他与他的祖先们一起埋葬在\Place{赫贝龙}。\par

% structure: p0069 source=h2 target=section reason=relative-heading-rank; base=h2

\section{V.— \Person{依撒加尔}的遗训：论纯朴}

1. 依撒加尔言论记录。他召集儿子们，对他们说：\Quote{我的孩子们，请听你们父亲\Person{依撒加尔}的话；你们这些蒙上主所爱的人啊，请侧耳听我的话。我是\Person{雅各伯}的第五个儿子，是因曼德拉草而得的工价。因为\Person{勒乌本}从田间带回曼德拉草，\Person{辣黑耳}遇见他，便拿走了。\Person{勒乌本}哭泣，他的哭声引来了我的母亲\Person{肋阿}。这些曼德拉草是一种气味甘甜的果实，产于\Place{阿兰}之地，生长在水沟旁的高地上。\Person{辣黑耳}说：\InnerQuote{我不给你，这些要代替儿女归我。}当时有两颗果子；\Person{肋阿}说：\InnerQuote{你夺去我丈夫的童贞还不够吗？还要夺走这些吗？}\Person{辣黑耳}说：\InnerQuote{看，今夜\Person{雅各伯}归于你，代替你儿子的曼德拉草。}\Person{肋阿}对她说：\InnerQuote{不要夸耀，不要自大；\Person{雅各伯}是我的，我是他青年时代的妻子。}\Person{辣黑耳}却说：\InnerQuote{怎么这样说？他先是许配给我的，并且为我缘故服事我们父亲十四年。我怎么办呢？因为人的诡诈和机巧增多了，诡诈在地上兴盛。若非如此，你现在不会看见\Person{雅各伯}的面。你不是他的妻子，而是藉诡计代替我嫁给他。我父亲欺骗了我，在那夜将我挪开，不让我见他；我若在那里，这事就不会发生。}\Person{辣黑耳}说：\InnerQuote{你取一个曼德拉草，另一个你夜间从我租用他。}\Person{雅各伯}亲近\Person{肋阿}，她怀孕生了我，因这工价，我名叫\Person{依撒加尔}。}\par

2. 那时，上主的一位天使向\Person{雅各伯}显现说：\Quote{辣黑耳要生两个孩子；因为她拒绝与丈夫同房，选择了节欲。若不是我母亲\Person{肋阿}为了同房而交出那两个苹果，她本可生八个儿子；但因这事她生了六个，\Person{辣黑耳}生了两个：因为上主因曼德拉草眷顾了她。祂知道她愿为儿女与\Person{雅各伯}同房，并非出于淫欲。她更进一步，次日又将\Person{雅各伯}让出，为要得另一个曼德拉草。因此，上主因曼德拉草俯听了\Person{辣黑耳}：因为她虽想要，却没有吃，而是带给了当时的至高者司铎，献于上主的殿中。}\par

3. \Quote{因此，我的孩子们，当我长大以后，我心怀正直而行，成为我父母和兄弟的农夫，按季节从田间带回果实；我父亲祝福我，因他见我行事纯朴。我行事不多管闲事，也不怀恶意诽谤邻人。我从未攻击任何人，也不指责任何人的生活，而是以纯朴的目光行事。因此，当我三十岁时，我娶了妻子，因劳作耗尽了我的精力，我从未思念与女人的欢愉；但我的劳动足以使我安眠，我父亲常因我的纯朴而喜乐。因我无论劳作什么，都首先通过司铎之手，将一切出产和初熟之物献给上主；然后给我父亲，再自取所需。上主加倍使祂的恩惠临于我的双手；\Person{雅各伯}也知道天主助我纯朴，因我以纯朴之心，将地上的美物施予每个穷人和困苦之人。}\par

4. \Quote{如今你们要听我，我的孩子们，以纯朴之心行事，因我见其中一切皆蒙上主悦纳。纯朴者不贪求金子，不欺骗邻人，不渴求多样美食，不喜各种衣饰，不为自己设想长寿，只等待天主的旨意，错谬之灵不能胜过他。因他不能容许心中存有女性之美貌的念头，以免以败坏污染自己的心思。嫉妒不能进入他的意念，妒火不融化他的灵魂，他也不以无尽贪欲图谋利益；因他生活正直，以纯朴观看万物，不容眼中因世间的错谬而有恶意，以免他看见任何违反上主诫命的乖谬。}\par

5. \Quote{因此，我的孩子们，要遵守天主的法律，获得纯朴，行事纯真，不可过分好奇地探究天主的诫命和邻人的事务；但要爱上主和你的邻人，怜恤穷人和软弱者。屈身务农，以各种农耕方式耕种土地，以感恩之心向上主献上礼物；因上主以地上的初熟之物祝福了我，一如祂从\Person{亚伯尔}直到如今祝福了所有圣人。因为除了土地的肥腴，没有别的份分给你们，其果实由劳作而得；因我父\Person{雅各伯}以地上的祝福和初熟之物祝福了我。\Person{肋未}和\Person{犹大}在\Person{雅各伯}的子孙中蒙上主荣耀；因上主拣选了他们，赐给一位司铎职，赐给另一位王国。因此你们要服从他们，并照你们父亲的纯朴行事；因为\Person{加得}被赐予摧毁那将要临于\Place{以色列}的试探。}\par

6. \Quote{我的孩子们，我知道在末后的日子，你们的子孙将离弃纯朴，依附贪婪，离弃纯真而接近恶毒，离弃上主的诫命而依附\Person{贝里雅耳}，离弃农耕而随从其邪恶的计谋，并被分散到外邦人中，服事他们的仇敌。因此你们要将这些事吩咐你们的子孙，使他们若犯罪，能更快地归向上主；因为祂是仁慈的，必拯救他们，甚至把他们带回他们的土地。}\par

我现年一百二十二岁，自知并未犯过至于死的罪。除了我的妻子，我未曾认识任何女人；我从未因眼目高傲而行淫；我没有饮酒以至于被引入歧途；我没有贪恋邻舍任何可欲之物；诡诈从未进入我心；谎言从未经过我口；若有人忧伤，我与他一同哭泣，并与穷人分享我的饼。我从未独自进食；我没有挪移地界；在我所有的日子中，我行了虔诚与真理。我全心爱了上主；同样，我也爱每个人如同自己的儿女。所以，我的孩子们，你们也要做这些事，那么贝利亚的一切邪灵必从你们面前逃跑，恶人的任何行为必不能管辖你们；你们要制伏一切野兽，因为你们拥有与心怀纯朴之人同行的天主。\par

于是他吩咐他们将他抬到\Place{赫贝龙}，与他的祖先一同葬在那里的洞穴中。他伸了双脚，安然去世，雅各伯的第五个儿子，寿高年迈，肢体健全，精力未衰，睡了永眠。\par

% structure: p0078 source=h2 target=section reason=relative-heading-rank; base=h2

\section{VI.— 则步隆的遗训：论怜悯与慈悲。}

1. 则步隆的记述，即他在一百一十四岁、若瑟去世三十二年后所嘱咐其子女的。他对他们说：\Quote{则步隆的儿子们，请听我言，听从你们父亲的话。我是则步隆，是我父母的美好恩赐。因为我出生时，我们的父亲因使用条纹木棍而分得福分，在羊群和牛群上极其丰盛。我的孩子们，我不知道在我所有的日子中，我是否犯了罪，除非仅在思想上。我也不记得我做过任何不义，除了我对若瑟所犯的无知之罪；因为我遮掩了我的兄弟们，没有告诉父亲所发生的事。我暗中痛哭，因为我惧怕我的兄弟们，因为他们都约定，若有人泄露秘密，必被刀剑所杀。但当他们想要杀他时，我含泪恳求他们不要犯下这罪。}\par

2. 因为西默盎和加得前来要杀若瑟。若瑟俯伏在地，对他们说：\Quote{我的兄弟们，可怜我吧，怜悯我们父亲雅各伯的心肠吧：不要伸手流无辜之血，因为我并未得罪你们；即使我得罪了，你们就责打我，但不要为我父亲雅各伯的缘故伸手。}他说这些话时，我怜悯他，开始哭泣，我的心在我里面融化了，我五脏六腑都松软了。若瑟也哭了，我也与他同哭；我的心剧烈跳动，我身体的关节颤抖，我站立不住。当他看见我与他同哭，而他们前来要杀他时，他逃到我身后，恳求他们。勒乌本说：\Quote{我的兄弟们，我们不要杀他，不如把他丢进我们祖先所挖、无水可寻的一个干坑里。因为上主为此禁止水在其中涌出，以便若瑟得以保全；上主如此安排，直到他们把他卖给依市玛耳人。}\par

3. 因为就若瑟的价钱而言，我的孩子们，我没有分得；但西默盎、加得以及我们另外六个兄弟取了若瑟的价钱，为自己、妻子和儿女买了凉鞋，说：\Quote{我们不吃这钱，因为它是我们兄弟的血价，我们要把它踩在脚下，因为他曾说他是我们的王，让我们看看他的梦是什么意思。}因此，在\OriginalTerm{哈诺客}{Enoch}的法律书卷中记载：凡不愿为其兄弟立嗣者，要脱去他的凉鞋，并要吐唾沫在他脸上。若瑟的兄弟们不愿他们的兄弟存活，上主就解开了若瑟的凉鞋。因为他们来到埃及时，在门口被若瑟的仆人解开了凉鞋，遂按法郎的样式向若瑟下拜。他们不仅向他下拜，而且立即仆倒在他面前，还受人吐唾沫，因此他们在埃及人面前蒙羞；因为此后埃及人听到了我们所行的对若瑟的一切恶事。\par

4. 这些事之后，他们拿出食物；而我因怜悯若瑟，两天两夜没有尝过任何东西。犹大没有与他们一同进食，而是看守着坑；因为他怕西默盎和加得跑回去杀他。当他们看见我也不吃时，就派我去看守他，直到他被卖。他在坑里待了三天三夜，最后饥饿地被卖掉了。勒乌本听说他不在时若瑟已被卖，就撕裂衣服哀悼说：\Quote{我怎能面对我父亲雅各伯的面？}他拿起钱，追赶商人，却没有找到任何人；因为他们已经离开大路，匆忙穿过崎岖的小道。勒乌本当日没有进食。于是丹来到他跟前说：\Quote{不要哭，也不要悲伤；因为我找到了我们可以对我们父亲雅各伯说的话。我们杀一只山羊羔，把若瑟的外衣蘸在里面；然后\InnerQuote{看看这是不是你儿子的外衣？}因为他们正要卖若瑟时，从他身上剥下了我们父亲的外衣，给他穿上了一件旧奴隶的衣服。}西默盎拿着那外衣，不肯交出，想用剑把它撕碎；因为他恼怒若瑟还活着，自己未能杀死他。于是我们全体起来反对他，说：\Quote{你若不肯交出，我们就说唯独你在以色列行了这恶事；}于是他就交出来了，他们便照丹所说的话做了。\par

5. 如今，我的孩子们，我吩咐你们遵守主的诫命，对邻人施行仁慈，对一切怀有怜悯，不仅对人，也对走兽。因为为此缘故，主祝福了我；当我所有的弟兄都患病时，我得以无病无恙，因为主知道各人的意图。因此，我的孩子们，要在你们心中怀有怜悯，因为人怎样对待邻人，主也必照样对待他。我弟兄的儿子们因若瑟的缘故生病、死亡，因为他们心中没有显出仁慈；但我的儿子们得以无病无恙，如你们所知。当我在客纳罕沿海之地时，我为我的父亲雅各伯捕获鱼获；当许多人在海中窒息时，我安然无恙。\par

6. 我是第一个制造船只在海上航行的人，因为主赐给我对此的理解和智慧；我在船后放下舵，在中间竖立一根桅杆，上面张起帆；乘船沿岸航行，我为父亲的家捕鱼，直到我们进入埃及；出于怜悯，我将我的鱼分给每个陌生人。若有人是陌生人、患病或年迈，我就煮鱼，烹调美味，供给众人各取所需，将他们聚集起来，怜悯他们。因此主也赐我捕获许多鱼：因为凡分施给邻人的，必从主获得更多。我捕鱼五年，将鱼分给我所见的每个人，并足以供给父亲全家。夏天我捕鱼，冬天我与弟兄们一同牧羊。\par

7. 现在我要向你们宣告我所行的事：我在冬天看见一个人困苦赤身，便怜悯他，偷偷从家中取出一件衣服，给了那困苦之人。因此，我的孩子们，你们要因天主所赐予你们的，向众人公平地施行怜悯和仁慈，以善良的心给予每个人。若你们当时没有东西给那求告你们的人，也要以怜悯的肠肺对他怀有同情。我知道我的手当时没有东西给那求告我的人，我便与他同行，哭泣超过七斯塔狄，我的肠肺向他涌动着怜悯。\par

8. 因此，我的孩子们，你们自己也要以仁慈向每人大发怜悯，好使主也怜悯你们，向你们施行仁慈；因为在末世，天主将祂的怜悯降在地上，凡祂找到有怜悯肠肺之处，祂便居住在其中。因为人怎样怜悯邻人，主也照样怜悯他。当我们下到埃及时，若瑟没有对我们怀有恶意，当他看见我时，他充满了怜悯。我的孩子们，你们也要效法他，证明自己毫无恶意，彼此相爱；不要各人计较弟兄的恶，因为这破坏合一，分裂一切亲族，扰乱灵魂：因为怀有恶意的人没有怜悯的肠肺。\par

9. 请看那水流：它们汇合在一起，冲走石头、树木、沙土；但若它们分成许多支流，土地就会吸干它们，它们便成为无用的。同样，你们若分裂，也将如此。不要分成两个头，因为主所造的一切只有一个头；祂赐下两个肩膀、双手、双脚，但所有肢体都服从于一个头。我从我祖先的著作中得知，在末世你们将离弃主，在以色列中分裂，跟随两个君王，行一切可憎的事，敬拜一切偶像，你们的仇敌将掳掠你们，你们将在万民中居住，充满各种疾病、困苦和灵魂的痛苦。此后你们将记起主，并悔改，祂将带领你们归回；因为祂是仁慈的，充满怜悯，不将恶归于人子，因为他们是血肉之躯，而谬误之灵在他们一切作为中欺骗他们。此后，主亲自向你们兴起，\BibleRef{拉 4:2}，正义之光，医治和怜悯在祂的翼上。祂要从贝里雅耳救赎人子的一切被掳，一切谬误之灵要被践踏。祂要使万民归向对祂的热忱，你们将在主所拣选的人形中看见天主，耶路撒冷是祂的名。你们又将因你们言语的邪恶激怒祂，你们将被抛弃，直到终结之时。\par

10. 如今，我的孩子们，不要因我死亡而悲伤，也不要因我离开你们而烦恼。因为我将再次在你们中间复活，如同一个统治者在他儿子们中间；我将在我的支派中欢欣，就是那些遵守了天主的法律和他们的父亲则步隆的诫命的人。但主将把永火降在不虔敬者身上，世世代代消灭他们。我正急忙前往我的安息之所，如同我的祖先一样；但你们要一生尽力敬畏你们的天主。他说完这些话后，平静地睡着了，他的儿子们将他放在棺木中；后来他们把他运到赫贝龙，与他的祖先同葬。\par

% structure: p0089 source=h2 target=section reason=relative-heading-rank; base=h2

\section{VII. — 关于愤怒与谎言之丹的遗嘱。}

1. 以下是\Person{丹}的遗言记录，即他在末日向儿子们所说的话。在他生命的第125年，他召集全家说：\Quote{我儿丹，请听我的话语；留心听你们父亲口中的言语。我在心中并在我整个生命中已证实：真理与公正的行为是善的，且中悦天主；而谎言与忿怒是恶的，因为它们教人一切邪恶。今天我向你们，我的孩子们，承认：我心中曾为真善人\Person{若瑟}的死而欢喜；我也为\Person{若瑟}被卖而欢喜，因为他的父亲爱他胜过爱我们。因为嫉妒和虚荣之灵对我说：\InnerQuote{你也是他的儿子。}\Person{贝里雅尔}的一个邪灵与我同工，说：\InnerQuote{拿起这把剑，用它杀死\Person{若瑟}；这样，他死了，你父亲就会爱你。}这是忿怒之灵所劝我的，叫我如豹吞食羔羊般吞食\Person{若瑟}。但我们祖先\Person{雅各伯}的天主没有把他交在我手中，叫我单独遇见他，也不容许我犯这罪，以免\Place{以色列}的两个支派被毁灭。}\par

\Quote{如今，我的孩子们，我快要死了，我如实告诉你们：除非你们谨守自己远离谎言和忿怒之灵，并爱慕真理和忍耐，你们必会灭亡。忿怒中有盲目，我的孩子们，没有哪个发怒的人能按真理看待任何人：因为即使那是父亲或母亲，他也待之如仇敌；即使是兄弟，他也不认识；即使是主的先知，他也不听从；即使是义人，他也不在意；朋友，他也不承认。因为忿怒之灵用欺骗的网罗围绕他，蒙蔽他本性的眼睛，藉着谎言暗昧他的心智，给他一个自我制造的幻象。他用什么包围他的眼睛呢？是用心中的憎恨；他给他一颗敌对兄弟直至嫉妒的心。}\par

\Quote{我的孩子们，忿怒是有害的，因为它对灵魂本身而言如同另一个灵魂；它使发怒之人的身体成为它的，并获取对灵魂的统治权，将其自身的力量赋予身体，以行一切不义；每当灵魂作什么事，它便为所作的事辩护，因为它看不见。因此，发怒的人若是强者，他发怒时有三种力量：一是藉仆人的力量和援助，二是藉他的忿怒，以此说服并靠不义取胜；三是他自身身体的本性，以及他自己作恶的本性。即使发怒的人是弱者，他也有比本性力量加倍的力量；因为忿怒常助长此类人的恶行。这灵常与谎言一同伴随在\Person{撒殚}的右边，好使他的作为得以藉残忍和谎言实现。}\par

\Quote{因此你们要明白忿怒的力量是虚妄的。因为首先它用言语刺痛人；然后藉行为强化发怒的人，并用痛苦的惩罚搅乱他的心智，从而以极大的忿怒激动他的灵魂。所以，当有人攻击你们时，你们不要被激动而发怒。若有人称赞你们为善，也不要自高自大或得意，无论是内心的感受还是外表的喜悦。因为首先它取悦听觉，然后激起理解力去领会忿怒的缘由；随后，发怒的人便以为自己是理直气壮地发怒。我的孩子们，若你们遭遇任何损失或毁灭，不要烦恼；因为这灵使人渴望已失去之物，好让他们被欲望点燃。若你们自愿承受损失，不要烦恼，因为烦恼会藉谎言引发忿怒。而忿怒与谎言是双重的祸害；它们互相勾结，以搅乱心智；当灵魂不断被搅乱时，主便离开它，\Person{贝里雅尔}便统治它。}\par

\Quote{因此，我的孩子们，要遵守主的诫命，持守祂的法律；远离忿怒，憎恶谎言，好使主住在你们中间，\Person{贝里雅尔}从你们面前逃遁。你们各人要与邻居说实话，这样你们就不会陷入情欲和混乱；但你们要在平安中，拥有和平的天主，这样就没有战争能胜过你们。你们一生要爱主，并彼此以真诚的心相爱。因为我知道在末日你们将离弃主，激怒\Person{肋未}，与\Person{犹大}争战；但你们不能胜过他们。因为主的天使必引导他们二人；因着他们\Place{以色列}才得站立。当你们离弃主时，你们将行一切邪恶，做外邦人的可憎之事，与不虔敬者的女子同流合污；谬误之灵将用一切恶念在你们里面作工。因为我在\WorkTitle{义人哈诺客书}中读到：你们的君王是\Person{撒殚}，一切邪淫和骄傲的灵都将服从\Person{肋未}，为给\Person{肋未}子孙设下陷阱，使他们在主面前犯罪。我的儿子们将亲近\Person{肋未}，在一切事上与他们一同犯罪；\Person{犹大}的儿子们将贪婪，像狮子一样抢夺别人的财物。因此你们将与他们一同被掳去，在那里你们将遭受\Place{埃及}的一切灾殃和外邦人的一切恶毒；然后，当你们归向主时，你们必得怜悯，祂必带你们进入祂的圣所，向你们宣告和平；从\Person{犹大}和\Person{肋未}支派中，必有主的救恩兴起；祂必与\Person{贝里雅尔}争战，并将胜利的报复赐给我们的疆界。祂必从\Person{贝里雅尔}手中夺回被掳的，即圣者的灵魂，并使悖逆的心归向主，并将永远的平安赐给呼求祂的人；圣者将在\Place{伊甸园}安息，义人将在\Place{新耶路撒冷}喜乐，这城将归于天主的荣耀，直到永远。\Place{耶路撒冷}不再荒凉，\Place{以色列}不再被掳；因为主必在她中间，住在人中间，即\Place{以色列}的圣者以谦卑和贫穷统治他们；信祂的人必在真理中统治诸天。}\par

6. 如今，我的孩子们，要敬畏主，并谨防撒殚及其邪魔；要亲近天主，和那为你们代祷的天使，因为祂是调停天主与人之间和平的天主与人的中保。祂必要起来对抗敌人的王国；因此，敌人切望摧毁所有呼求主名的人。因为他知道，在以色列相信的那一天，敌人的王国必要完结；那和平的天使必要坚固以色列，使它不至陷入极恶之中。在以色列的邪恶时期，主必离开他们，而跟随那承行祂旨意的人，因为祂的众天使中，没有一个能像祂那样。祂的名在以色列各处及外邦人中，必称为救主。所以，我的孩子们，要保守自己远离一切恶行，弃绝忿怒和一切谎言，喜爱真理和忍耐；你们从你们父亲所听到的，也要传授给你们的子女，好使外邦人的父接纳你们：因为祂是真实而忍耐的，温良而谦卑，并以祂的作为教导天主的法律。因此，要远离一切不义，持守主的法律之正义：并将我埋葬在我祖先的身旁。\par

7. 他说完这些话后，就亲吻了他们，然后长眠了。他的儿子们埋葬了他，随后将他的骸骨运到亚巴郎、依撒格和雅各伯的旁边。然而，正如丹曾预言他们将忘记他们天主的法律，并与他们继承的土地、以色列的种族和他们的亲属疏远，事情果然如此发生了。\par

% structure: p0097 source=h2 target=section reason=relative-heading-rank; base=h2

\section{八、—— 纳斐塔里关于自然良善的遗训。}

1. 纳斐塔里遗训的记录，即他在世第一百三十二年临终时所吩咐的事。当他的儿子们在第七个月、第四日聚集时，他尚健在，为他们设宴欢庆。早晨醒来后，他对他们说：我要死了；他们却不肯相信。他祝福了主；并断言在昨天的宴席之后他要死去。于是他开始对他的儿子们说：我的孩子们，请听；你们纳斐塔里的儿子们，请听你们父亲的话。我是从彼耳哈所生；因为辣黑耳用了计谋，将彼耳哈代替自己给了雅各伯，她在我出生时使我落在辣黑耳的膝上，因此我被称为纳斐塔里。辣黑耳爱我，因为我是在她膝上出生的；当我年幼体弱时，她常亲吻我说：但愿我能在我的腹中见到一个像你一样的兄弟；因此，因着辣黑耳的祈祷，若瑟在各方面都像我。我的母亲是彼耳哈，她是勒贝加乳母德波辣的兄弟罗特乌斯的女儿，她与辣黑耳同一天出生。罗特乌斯是亚巴郎家族的人，是加色丁人，敬畏天主，生而自由且高贵；他被掳后，被拉班买下；拉班将他的婢女爱纳给他为妻，她生了一个女儿，取名齐耳帕，以他被掳村庄的名字命名。接着她又生了彼耳哈，说：我的女儿渴望新奇事，因为她一出生就急切要吃奶。\par

2. 由于我奔跑如鹿般迅捷，我的父亲雅各伯派我处理一切差事和信使任务，并像一只鹿（\BibleRef{创 49:21}）那样赐福了我。因为正如陶匠知道器皿的用途，知道它盛装什么，并为之加上泥土，同样，主也依照精神塑造身体，并依照身体的能力植入精神，两者之间连一根头发的三分之一也不欠缺；因为至高者的每个受造物都是按重量、尺度和规则造成的。正如陶匠知道每个器皿的用途，知道它适宜做什么，同样，主也知道身体，知道它在善事上能到何种程度，何时开始作恶；因为没有一样受造物或思想是主所不知道的，因为祂按自己的形象造了每个人。一个人的力量如何，他的工作也如何；他的心意如何，他的工作也如何；他的目的如何，他的作为也如何；他的心如何，他的口也如何；他的眼如何，他的睡眠也如何；他的灵魂如何，他的言语也如何，要么在主的法律中，要么在贝里雅耳的法律中。正如光与暗、看与听之间有区别，人与人、女人与女人之间也有区别；也不能说任何事物有优越性，无论是面貌还是其他类似之物。因为天主按秩序造了一切美好之物：头部的五种感官，祂将颈项接在头上，头发也为美观，心为悟性，腹为胃的分化，芦管为健康，肝为怒气，胆为苦味，脾为欢笑，肾为机巧，腰为力量，肋为容纳，背为强壮，等等。所以，我的孩子们，要在敬畏天主中按秩序行善，不要因轻蔑或不合时宜而做无序之事。因为如果你叫眼睛去听，它不能；同样，在黑暗中你也无法行光明的事。\par

3. 因此，不要因过度而急于败坏你们的行为，或以虚言欺骗你们的灵魂；因为如果你们以纯洁的心保持缄默，就能持守天主的旨意，弃绝魔鬼的意志。太阳、月亮和星辰不改变它们的秩序；同样，你们也不要以行为的无序改变天主的法律。万民迷失了，离弃了主，改变了他们的秩序，跟随石头和木头，追随错谬的神灵。但你们不可如此，我的孩子们，要认识那在穹苍、大地、海洋和一切受造物中创造万有的主，免得你们像索多玛人那样改变了本性的秩序，同样，守夜者也改变了他们本性的秩序，主在洪水时诅咒了他们，并为了他们的缘故使大地荒废，无人居住，不结果实。\par

4. 这些事，我的孩子们，我告诉你们，因为我在\Person{以诺}的神圣著作中读到，你们自己也将离弃主，随从外邦人的一切恶行，并要行索多玛的一切不义。主必将掳掠临到你们，你们要在那里服侍你们的仇敌，被各种苦难和患难所覆盖，直到主将你们全部消灭。此后，你们人数减少，所剩不多，你们将回转，承认主你们的天主；他必照他丰盛的怜悯，将你们带回你们自己的土地。并且，在他们进入他们祖先的土地之后，他们将再次忘记主，行恶；主必将他们分散在全地面上，直到主的怜悯降临，一位行义并施仁慈的人\OriginalTerm{人}{man}，向远近所有的人施恩。\par

5. 因为在我四十岁那年，我\Emph{在异象中}看见太阳和月亮停在橄榄山上，在耶路撒冷以东。看哪，我父亲的父亲\Person{依撒格}对我们说：\Quote{跑去，各按自己的力量抓住它们；谁抓住了，太阳和月亮就归谁。}我们全都一起跑，\Person{肋未}抓住了太阳，\Person{犹大}超过了其他人抓住了月亮，他们两个都被举起。当\Person{肋未}变得像太阳时，一个年轻人给了他十二枝棕榈枝；\Person{犹大}明亮如月亮，在他脚下有十二道光芒。\Person{肋未}和\Person{犹大}奔跑，彼此抓住。看哪，地上有一头公牛，长着两只大角，背上有一对鹰翅；我们想抓住它，却不能。因为\Person{若瑟}超过了我们，抓住了它，并带着它升到高处。我看见——因为我就在那里——看哪，一篇神圣的著作出现在我们面前，说：\Quote{亚述人、玛待人、波斯人、厄蓝人、革拉基人、加色丁人、叙利亚人，将掳掠以色列十二支派。}\par

6. 又过了七个月，我看见我们的父亲\Person{雅各伯}站在雅木尼雅海边，我们他的儿子们与他在一起。看哪，有一艘船航行而来，满载干肉，没有舵手或领航员；船上写着：\Quote{\Person{雅各伯}。}我们的父亲对我们说：\Quote{我们上我们的船吧。}我们上船后，起了猛烈的风暴，强风大作；我们的父亲握着舵，却从我们面前飞走了。我们被风暴抛掷，漂流海上；船里满了水，被大浪拍打，几乎碎裂。\Person{若瑟}驾着一只小船逃走了，我们全部分散在十二块木板上，\Person{肋未}和\Person{犹大}在一起。于是我们全都四散到远方。然后，\Person{肋未}腰束麻布，为我们众人向主祈祷。风暴止息时，船立刻靠岸，如同平安一样。看哪，我们的父亲\Person{雅各伯}来了，我们同心欢喜。\par

7. 这两个梦我告诉了我的父亲；他对我说：\Quote{这些事必在它们的时期应验，在以色列经历许多事以后。}然后我的父亲对我说：\Quote{我相信\Person{若瑟}还活着，因为我常看见主把他列在你们中间。}他哭着说：\Quote{你活着，我的孩子\Person{若瑟}，我却看不见你，你也看不见生你的\Person{雅各伯}。}他的话也使我们哭泣；我心里发热，想要说出他被卖的事，但我害怕我的弟兄们。\par

8. 看哪，我的孩子们，我已将末后的日子指示给你们，这一切都将发生在以色列。因此，你们也要嘱咐你们的儿女，叫他们与\Person{肋未}和\Person{犹大}联合。因为救恩必从\Person{犹大}临到以色列，在他内\Person{雅各伯}必蒙祝福。因为通过他的支派，天主要被看见在地上居住在人间，以拯救以色列的种族，并要从外邦人中聚集义人。我的孩子们，如果你们行善，人和天使都要祝福你们；天主也要因你们在外邦人中得荣耀，魔鬼要从你们面前逃跑，野兽要惧怕你们，天使要依附你们。因为正如一个人好好养育一个孩子，对他有美好的纪念；同样，善行在天主面前也有美好的纪念。但那不行善的人，人和天使都要咒诅他，天主也要因他在外邦人中受羞辱，魔鬼使他成为自己特有的工具，一切野兽都要制服他，主也要恨他。因为法律的诫命是双重的，必须通过智慧来履行。因为人拥抱妻子有时节，为祈祷而禁欲也有时节。所以有两条诫命；若不按适当的次序实行，就带来罪。其他的诫命也是如此。因此，你们要在天主内明智，有智慧，明白诫命的次序和每项工作的法则，好使主爱你们。\par

9. 他用许多这样的话嘱咐了他们之后，便劝勉他们把他的骸骨迁到赫贝龙，与他祖先同葬。他欢欢喜喜地吃喝之后，蒙上脸，就死了。他的儿子们完全遵行了他们的父亲\Person{纳斐塔里}所嘱咐他们的一切事。\par

% structure: p0107 source=h2 target=section reason=relative-heading-rank; base=h2

\section{九．——\WorkTitle{加得遗训}：论仇恨。}

1. 加得的遗训记录，即他在一百二十七岁时向众子所说的话：\Quote{我是雅各伯所生的第七子，在守护羊群上勇敢。我夜间看守羊群；每当狮子、豺狼、豹子、熊或任何野兽前来袭击羊圈，我便追赶它，用手抓住它的脚，旋转它，使它晕眩，然后抛掷到两弗隆之外，就这样杀死了它。那时若瑟与我们一同牧羊约三十天，因年幼体弱，又因酷热而病倒。他回到赫贝龙他父亲那里，父亲因爱他，便让他睡在自己身边。若瑟告诉父亲说，齐耳帕和彼耳哈的儿子们宰杀最好的牲畜，背着犹大和勒乌本偷吃。原来他看见我从熊口中救出一只羔羊，我杀死了熊；但那只羔羊因受伤无法存活，我便杀了它，我们吃了它，而他把这事告诉了父亲。因此我向若瑟恼怒，直到他被卖到埃及之日。仇恨的灵在我里面，我不愿看见若瑟，也不愿听他说话。他当面指责我们未经犹大同意就吃了羊群中的肉。凡他告诉我父亲的事，父亲都相信他。}\par

2. \Quote{我的孩子们，我现在承认我的罪过：我多次想要杀死他，因为我恨他入骨，对他毫无怜悯之心。此外，因他的梦，我更加恨他；我恨不得将他从活人之地吞灭，如同牛犊从地上吞吃青草一般。因此，我和犹大以三十块金子将他卖给依市玛耳人，其中十块我们藏了起来，只拿二十块给兄弟们看；这样，因我的贪心，我全心图谋他的毁灭。但我祖先的天主从他手中救出了我，使我在以色列中不致行恶。}\par

3. \Quote{现在，我的孩子们，请听真理之言，以行正义，并遵守至高者的一切法律，不要因仇恨之灵而误入歧途，因为它在人的一切行为中都是邪恶的。无论一个人做什么，仇恨者都厌恶他：即使他遵行主的法律，仇恨者也不称赞他；即使他敬畏主，喜爱正义的事，仇恨者也不爱他：他诽谤真理，嫉妒行事正直的人，喜爱恶言，傲慢自大，因为仇恨蒙蔽了他的灵魂——正如我当初看待若瑟一样。}\par

4. \Quote{因此，我的孩子们，要提防仇恨，因为它对主自身行不义：它不愿听祂关于爱近人的诫命，反而得罪天主。因为若一个弟兄失足，仇恨者立刻希望向众人宣扬，并急于使他因此受审判、受罚、被杀。若是一个仆人，他便向主人控告他，用尽一切折磨算计他，若有可能就杀死他。因为仇恨藉嫉妒工作，嫉妒那些在善行上顺利的人，一看见或听见便嫉妒。因为正如爱能使死人复活，召回被判死的人，同样仇恨会杀死活人，连那些犯小过的人也不容他们存活。因为仇恨之灵与撒殚一同藉急躁之灵在一切事上致人死亡；但爱之灵与天主的法律一同藉坚忍使人得救恩。}\par

5. \Quote{仇恨是邪恶的，因它常与谎言为伴，反对真理；它使小事变大，将黑暗视同光明，称甘甜为苦，教唆诽谤、争斗、强暴及一切邪恶的过分行为；并以魔鬼的毒液充满人心。我的孩子们，我凭经验向你们说这些话，为使你们逃避仇恨，持守主的爱。正义驱逐仇恨，谦卑消灭仇恨。因为公正而谦卑的人羞于行恶，不是受别人责备，而是受自己内心的责备，因为主鉴察他的意念：他不毁谤任何人，因为对至高者的敬畏胜过仇恨。因害怕得罪主，他不愿对人行任何不义，连在思想中也不。这些事我是在为若瑟悔改之后才学到的。因为合乎敬虔的真悔改消灭不信，驱散黑暗，照亮双眼，赐知识给灵魂，引导心思走向救恩；那些从人那里未学到的，藉悔改便知道了。因为天主使我的心患病；若非我父亲雅各伯的祈祷干预，我的灵魂几乎要离去。因为人因何事犯罪，也因何事受罚。因我的心对若瑟毫无怜悯，我的心也照样受无怜悯之苦，受审判十一个月，正如我嫉妒若瑟直到他被卖的那段时日。}\par

6. \Quote{现在，我的孩子们，你们要彼此相爱，从心中除去仇恨，在行为、言语和心灵的意念上彼此相爱。因为在我父亲面前，我与若瑟说话和平；但我一出去，仇恨之灵便蒙蔽我的心，激动我的灵魂要杀他。所以你们要从心里彼此相爱；若有人得罪你，要温和地告诉他，驱散仇恨的毒液，不要心怀诡诈。若他承认并悔改，就宽恕他；若他否认，不可与他争论，以免他起誓，而使你双倍犯罪。争论时不要让外人听见你的秘密，免得他恨你而成为你的敌人，并大大得罪你；因为他常会以诡诈与你交谈，或以恶计胜过你，从自身汲取毒液。因此，若他否认，又被指证受辱而默不作声，不可继续试探他。因为否认者若悔改，便不再对你有过犯；而且他会尊敬你，惧怕你，与你和平相处。但若他无耻，坚持作恶，即使那时也要从心里宽恕他，将报应交给天主。}\par

7. 若有人比你更发达，不要忧伤，反要为他祈祷，使他得享完满的发达。因为或许这样对你有益；纵使他更被高举，也不要嫉妒，要记念所有血肉之躯都将死去；并要赞美天主，祂赐予众人美好有益之物。要寻求主的判断，这样你的心便得安息与平安。纵然有人以邪恶手段致富，如同我父亲的兄弟厄撒乌，也不要嫉妒；但要等候主的终结。因为祂或从恶人那里收回祂的恩惠，或仍留给悔改者，或为不悔改者永留惩罚。因为那免于嫉妒的穷人，凡事感谢主，在众人中是富足的，因为他没有对人心存邪恶的嫉妒。所以，要从你们灵魂中除去仇恨，以正直的心彼此相爱。\par

8. 你们也要将这些事告诉你们的子女，使他们尊敬犹大和肋未，因为主要从他们之中兴起一位拯救者给以色列。因为我知道，最后你们的子女要离弃他们，在主面前行尽一切邪恶、祸害与败坏。他稍事休息后，又对他们说：我的孩子们，要听从你们的父亲，将我葬在我祖先旁边。于是他收拢双脚，平安地睡了。五年之后，他们将他运上去，安放在赫贝龙，与他祖先同葬。\par

% structure: p0116 source=h2 target=section reason=relative-heading-rank; base=h2

\section{X. —— 阿协尔的遗嘱，论恶与善的两面。}

1. 阿协尔遗嘱的记录，即他在一百二十岁时对他儿子们所说的话。他尚健康时，对他们说：阿协尔的孩子们，要听从你们的父亲，我将向你们宣示天主眼中一切正直的事。天主赐给了人子两条道路、两种心思、两种行为、两个地方和两种结局。因此，万物都是成对的，彼此相应。有善与恶两条道路，我们胸中的两种心思辨别它们。所以，若灵魂以善为乐，它一切行为都在正义中；即使犯罪，也立即悔改。因为，他既以正义为念，又抛弃恶意，便立刻推翻邪恶，根除罪恶。但若他的心思转向邪恶，他一切行为就在恶毒中，他赶走善，接纳恶，受贝里雅耳统治；即使他行善，也把它扭曲为恶。因为每当他似乎开始行善时，却使行为的结局归于恶，因为魔鬼的宝藏充满了恶神的毒液。\par

2. 于是，他说，有一种灵魂为恶的缘故说善，其行为的结局导向祸害。有一个人，对那在恶中服侍他需要的人毫不怜悯；这事有两面，但整体是恶。又有一个人，爱那作恶的人；他同样住在恶中，因为他甚至愿为他的缘故在恶事上死去：关于这一点，显然有两面，但整体是恶行。即使有爱，也只是掩盖邪恶的恶毒，如同它有一个看似善的名号，但行为的结果趋向恶。另一个人偷窃、行不义、掠夺、欺诈，同时又怜悯穷人：这也有一面，但整体是恶。他欺诈邻人，激怒天主，对至高者发假誓，却又怜悯穷人；他轻蔑并激怒那颁赐律法的主，却使穷人得舒畅；他污秽灵魂，却使身体愉悦；他杀害多人，却怜悯少数人：这也有一面。另一个人犯奸淫和淫乱，又禁食；但他在禁食时行恶，凭其权势和财富使多人堕落，又因他极度的邪恶而遵行诫命：这也有一面，但整体是恶。这样的人有如猪或兔；因为它们一半洁净，但事实上却不洁净。因为天主在\WorkTitle{天上碑版}中已如此宣告。\par

3. 所以，我的孩子们，不要像他们那样有两副面孔，一面善一面恶；但只应持守善，因为天主安息于善，人也渴慕善。要逃离邪恶，用你们的善工摧毁魔鬼；因为那些有两副面孔的人不事奉天主，只事奉自己的私欲，好叫他们取悦贝里雅耳和像他们一样的人。\par

4. 因为善人，就是那些有一副面孔的人，即使被有两副面孔的人看作有错，在天主面前却是正义的。因为许多人在杀恶人时作了两样工，一恶一善；但整体是善，因为他根除并毁灭了那邪恶的。有人恨那施怜悯者，又亏待奸夫和贼：这也是一面两面，但整体是善的，因为他效法主的榜样，不把看似善的与真正恶的并列。另有人不愿与狂欢的人同享好日子，免得玷污自己的口和灵魂：这也有一面两面，但整体是善的，因为这样的人像公鹿和母鹿，在野性状态中似乎不洁，但它们全然洁净；因为他们怀着对天主的热情而行，并戒绝天主也憎恶且藉祂的诫命所禁止的事，他们从善中挡住恶。\par

5. 因此，我的孩子们，你们看，万物中都有两个，彼此对立，一个被另一个所隐藏。死亡接续生命，羞辱接续光荣，黑夜接续白昼，黑暗接续光明；万物都在白昼之下，正义的事在生命之下：因此永生也等待死亡。不可说真理是谎言，也不可说公义是错误；因为所有真理都在光明之下，如同万物都在天主之下。这一切我一生都经验过，没有偏离主的真理，我探究了至高者的诫命，以一副面孔、尽我全部力量走在善事中。\par

6. \Quote{所以，我的孩子们，你们也要遵守主的诫命，以纯一的面貌跟随真理，因为心怀二意的人要受双重的惩罚。要憎恨那与人为敌的谬误之灵。遵守主的法律，不要以恶为善；而要看那真正为善的事，并持守它，在主的全部诫命中，以祂为你们的交往，并在祂内安息：因为人所追求的目标显明他们的正义，并且要从撒殚的天使中认识主的天使。因为灵魂若在忧虑中离去，它必被那它在私欲和恶行中所事奉的那恶灵折磨；但如果灵魂安静而喜乐地认识了和平的天使，这天使将在生命中安慰他。}\par

7. \Quote{我的孩子们，不要像\Place{索多玛}那样，她不认识主的天使，以致永远灭亡。因为我知道你们将要犯罪，你们将被交在敌人手中，你们的土地将变得荒芜，你们将被分散到地极四方。你们在分散中将被视为无用的水，直到至高者眷顾大地；祂将作为人而来，与人一同吃喝，在平安中藉着水击碎龙的头。祂要拯救以色列和万民，天主在人的位格中说话。因此，要把这些事告诉你们的孩子们，使他们不悖逆祂。因为我在天上的书卷中读到，你们确实会悖逆祂，对祂行不虔之事，不遵从天主的法律，却遵从人的诫命。因此，你们将如我的兄弟\Person{加得}和\Person{丹}那样被分散，他们不认识自己的土地、支派和言语。但主将藉着信德，因祂仁慈怜爱的希望，为了\Person{亚巴郎}、\Person{依撒格}和\Person{雅各伯}的缘故，把你们聚集在一起。}\par

8. 他对他们说了这些话之后，便吩咐他们说：要把我葬在\Place{赫贝龙}。然后他安详地入睡，去世了；之后，他的儿子们照他所吩咐的，把他抬上去，与他的祖先同葬。\par

% structure: p0125 source=h2 target=section reason=relative-heading-rank; base=h2

\section{十一、\Person{若瑟}的遗嘱：论节制}

1. \Person{若瑟}的遗嘱记录。当他即将去世时，他召集了他的儿子们和弟兄们，对他们说：\Quote{我的孩子们和弟兄们，请听以色列所爱的\Person{若瑟}；我的儿子们，请听你们父亲的话。我一生中见过嫉妒和死亡，但我在主的真理中没有迷失。我的这些弟兄们恨我，但主爱了我；他们想杀我，但我祖先的天主保护了我；他们把我下到坑里，但至高者又把我提升上来；我被卖为奴隶，但主使我自由；我被掳去，但祂强有力的手扶持了我；我遭受饥饿，但主亲自养育了我；我孤身一人，但天主安慰了我；我患病，但至高者眷顾了我；我被囚在监里，但救主对我施恩；我被捆绑，但祂释放了我；我遭诽谤，但祂为我辩护；在埃及人的恶言中，祂拯救了我；在嫉妒和诡计中，祂高举了我。}\par

2. \Quote{这样，法老的膳长\Person{普提法尔}把他的家托付给我，我与一个无耻的女人争战，她催逼我与她同犯淫行；但我父亲\Person{以色列}的天主保护了我，使我免于烈火的焚烧。我被投入监狱，被鞭打，被嘲弄；但主使我在监狱长眼前蒙恩。因为主决不会弃绝那些敬畏祂的人，无论他们在黑暗中、在捆绑中、在患难中、还是在困乏中。因为天主不像人那样感觉羞耻，也不像世人那样惧怕，也不像生于尘土的人那样软弱，也不能被推开；但祂处处都在，以多种方式安慰人，暂时离去以试炼灵魂的意向。在十次试探中，祂证明我是可悦纳的，我在这一切中都忍耐了；因为忍耐是强大的法宝，而坚忍带来许多益处。}\par

3. \Quote{那埃及女人多少次以死亡恐吓我！多少次她把我交到刑罚之下，然后又叫我回来，当我拒绝与她交往时，她又恐吓我！她对我说：\InnerQuote{你要作我的主人，凡属我的都是你的，只要你把自己给我，你就如同我们的主人一样。}因此，我记起我父亲\Person{雅各伯}的祖辈的话，便进入我的内室，向主祈祷；在那七年中我守斋，我在我主人面前显得生活优裕，因为为天主守斋的人容貌美丽。若有人给我酒，我不喝；我守斋三天，把我的食物拿来给穷人和病人。我一早就寻求主，为那个\Place{门非斯}的埃及女人哭泣，因为她不停地烦扰我，夜间她以探访我为借口来找我；起初，因为她没有男儿，她假装把我当作儿子。我向主祈祷，她便生了一个男儿；因此有一段时间她把我当作儿子拥抱，而我却不知道。最后，她试图引诱我犯奸淫。当我觉察到时，我悲伤欲死；她出去后，我回过神来，为她哀哭了多日，因为我看到她的诡诈和欺骗。我向她宣告至高者的话，或许她会从她邪恶的情欲中回转。}\par

4. 她多少次用言语向我谄媚，如同对待一位圣人，言语中却带着诡诈，在她丈夫面前称赞我的贞洁，而当我们独处时却想要毁灭我。她公开赞颂我为贞洁之人，暗中却对我说：\Quote{不要怕我的丈夫；因为他已确信你的贞洁，即使有人向他告发我们，他也决不会相信。}我因此披上苦衣躺在地上，恳求天主，愿主救我脱离那\OriginalTerm{埃及人}{Egyptian}。她一无所成之后，又以求教为借口再次来找我，想要知道主的圣言。她对我说：\Quote{你若愿意我离弃我的偶像，就听从我；我将说服我的丈夫离弃他的偶像，我们将在你主的法律中行走。}我对她说：\Quote{主不愿敬畏祂的人留在不洁中，也不喜悦犯奸淫的人。}她沉默了，一心想要实现她的邪恶欲望。我更加致力于禁食和祈祷，求主救我脱离她。\par

5. 又在另一时候，她对我说：\Quote{你若不肯犯奸淫，我就杀死我的丈夫，这样我就可以合法地娶你为夫。}我听到这话，就撕裂了衣服，说：\Quote{妇人，敬畏主，不要行这恶事，免得你彻底毁灭；因为我要向众人宣告你的不虔之念。}她因此害怕，恳求我不要向任何人宣告她的邪恶。她离去了，用礼物安抚我，并送给我人子的一切欢乐。\par

6. 她派人送来洒了魔法的食物。当送食物的阉人来时，我举目看见一个可怕的人用盘子给我一把剑，我便知道她的计谋是要欺骗我的灵魂。他出去后，我哭了，没有尝那食物，也没有尝她任何其他食物。于是过了一天，她来察看食物，对我说：\Quote{这是怎么回事，你没有吃食物？}我对她说：\Quote{因为你用死亡充满它；你怎么说：\InnerQuote{我不接近偶像，只接近主}呢？所以现在要知道，我父亲的天主藉着一位天使向我启示了你的邪恶，我保留着它，为要定罪你，或许你会看见并悔改。但为使你知道，不敬虔之人的邪恶不能胜过那些在贞洁中敬畏天主的人，我就在她面前取了食物吃了，说：\InnerQuote{我祖先的天主和\OriginalTerm{亚巴郎的天使}{Angel of Abraham}与我同在。}}她俯伏在我脚前哭泣；我扶起她，劝诫她，她答应不再行这恶事。\par

7. 但因她的心专注于我，要行淫乱，她就叹息，面容沮丧。她丈夫见了她，便对她说：\Quote{你的面容为何沮丧？}她说：\Quote{我心中疼痛，我的灵呻吟压迫我。}他就安慰这个本无病痛的人。随后，趁丈夫还在外面，她冲进来对我说：\Quote{你若不同意我，我就上吊，或投井，或跳崖。}当我看见\OriginalTerm{贝利亚尔}{Beliar}的精神正困扰她时，我向主祈祷，并对她说：\Quote{你为什么烦乱不安，被罪恶蒙蔽？记住，你如果自杀，你丈夫的妾\OriginalTerm{塞通}{Sethon}·你的情敌，会打你的孩子，从地上除灭你的纪念。}她对我说：\Quote{看，你爱我；这对我已足够——你关心我的生命和我的孩子；我盼望我能满足我的欲望。}她不知道，我这样说是因为我的天主，而不是因为她。因为一个人若在邪恶欲望的激情前跌倒，他就成了奴隶，正如她一样。若他听到任何关于他所征服的激情的好话，他就接受它，用于他的邪恶欲望。\par

8. 我向你们宣示，我的孩子们，那时大约是第六时辰，她离开了我；我整日跪在主前，并持续整夜；将近黎明，我哭着起来，祈求脱离那\OriginalTerm{埃及人}{Egyptian}。最后，她抓住我的衣服，强行拉我同她交合。因此，当我看见她在疯狂中强行抓住我的衣服时，我赤裸着逃走了。她向丈夫诬告我，那埃及人将我下在监里，在他的家中；次日，鞭打了我之后，那埃及人把我送进他家里的监牢。因此，当我被锁链束缚时，那个埃及妇人因烦恼而病倒，她听我在黑暗的居所如何向主歌唱赞美，并以喜乐的声音欢喜，光荣唯独我的天主，因为藉着一个借口，我摆脱了那个埃及妇人。\par

9. 她多少次派人来对我说：\Quote{同意实现我的欲望，我就释放你脱离锁链，我使你脱离黑暗！}我甚至在思想上也没有倾向她。因为天主爱那在黑暗的洞穴中守贞禁食的人，胜过那在密室中无拘无束地奢华生活的人。凡生活在贞洁中，也渴望荣耀的人，若\OriginalTerm{至高者}{Most High}知道这对他有益，祂也把这荣耀赐给他，如同赐给我一样。多少次，尽管她生病，却在意想不到的时候下到我这里，听我祈祷的声音！我听到她的呻吟时，就保持沉默。因为我在她家里时，她惯常裸露她的手臂、胸部和腿，好使我屈服于她；因为她非常美丽，为欺骗我而打扮得光彩夺目。主保护了我，使我脱离她的诡计。\par

10. 因此，我的孩子，你们看，忍耐、祈祷与禁食成就了何等伟大的事。如果你们因此以忍耐和谦虚的心追求节制与纯洁，主就住在你们中间，因为祂喜爱节制。凡至高者居住的地方，即使一个人陷入嫉妒、为奴或诽谤，那住在他内的主，因他节制的缘故，不仅拯救他脱离邪恶，而且高举并光荣他，如同我一样。因为人无论行为、言语或思想，都受到完全的护卫。我的兄弟们知道我的父亲多么爱我，而我心中并没有自高；虽然我是孩童，我在思想中却敬畏天主。因为我知道万事都将过去，我持守本分，尊敬我的兄弟们；因着对他们的敬畏，我被卖时默默不语，没有向依市玛耳人透露我的家族，说我是雅各伯之子，一个伟大而有权势的人。\par

11. 因此，你们也要在行为中敬畏天主，并尊敬你们的弟兄。因为凡遵行主法律的人，必为主所爱。当我与依市玛耳人来到\Place{Indocolpitæ}时，他们问我，我便说我是他们家的奴仆，以免使我的兄弟们蒙羞。他们中最年长的对我说：\Quote{你不是奴仆，因为连你的外貌也证明了你不是。}他甚至以死威胁我。但我仍然说我是他们的奴仆。当我们进入埃及时，他们为了我而争执，谁该买下我并带走我。于是大家都认为，我应该留在埃及，与一个同行商人在一起，直到他们带着货物回来。主使我在那商人眼中蒙恩，他便把他的家托付给我。主因我而祝福他，使他的金银增多，我与他在一处共三个月零五天。\par

12. 大约那时，\Person{Potiphar}的孟斐斯妻子以极大的盛况经过，并注目于我，因为她的太监们向我提及了她。她便向她的丈夫提及那商人，说他藉着一个希伯来少年发了财，说：\Quote{而且他们说他确实是从客纳罕地被人偷来的。现在，你当审断他，带走这少年做你的管家；这样，希伯来人的天主必会祝福你，因为来自天上的恩宠在他身上。}\par

13. \Person{Potiphar}被妻子的话说服了，便下令将那商人带来，对他说：\Quote{我所听见的是什么？你从希伯来地偷窃灵魂，卖他们为奴？}那商人于是俯伏在地，恳求他说：\Quote{我恳求你，我的主，我不知道你在说什么。}他说：\Quote{那么你的希伯来仆人从何而来？}他说：\Quote{依市玛耳人将他托付给我，直到他们回来。}他不信他，却下令剥去他的衣服并鞭打他。当那商人坚持己见时，\Person{Potiphar}说：\Quote{把那少年带来。}当我被带进来时，我向太监长行礼——因为他在法郎面前位列第三，是众太监之长，并有妻子、儿女和妾侍。他把我拉到一旁，对我说：\Quote{你是奴仆还是自由人？}我说：\Quote{奴仆。}他对我说：\Quote{你是谁的奴仆？}我回答他说：\Quote{依市玛耳人的。}他又对我说：\Quote{你怎么成了他们的奴仆？}我说：\Quote{他们从客纳罕地买了我。}他不信我，说：\Quote{你在说谎。}并下令剥去我的衣服鞭打我。\par

14. 那孟斐斯妇人正透过窗户观看我受鞭打，便派人到她丈夫那里说：\Quote{你的审判不公；因为你甚至惩罚一个被偷来的自由人，好像他是犯法者。}尽管我受鞭打，却没有别的回答；于是他下令将我们看管起来，直到他说\Quote{这少年的主人会来}为止。他的妻子对他说：\Quote{你为什么囚禁这高贵的少年？他理当得到自由，并侍奉你。}原来她怀着罪恶的欲望想见我，而我对此一无所知。于是他对妻子说：\Quote{埃及人没有在证据确凿之前夺取他人之物的习俗。}他这样说是指那商人和我，认为我必被囚禁。\par

15. 二十四天后，依市玛耳人来了；他们听说我的父亲雅各伯因我而哀恸，便对我说：\Quote{你怎么说你是个奴仆？看，我们得知你是客纳罕地一个大能之人的儿子，你的父亲正披麻为你哀恸。}我又想哭泣，但我克制自己，以免使我的兄弟们蒙羞。我说：\Quote{我不知道，我是一个奴仆。}于是他们商议要把我卖掉，免得我落在他们手中。因为他们惧怕雅各伯，怕他向他们施以致命的报复。因为听说他与主及与人都大有能力。那商人对他们说：\Quote{请使我脱离\Person{Potiphar}的审判。}他们便来要我说：\Quote{他是我们用钱买来的。}他便把我们放了。\par

16. 那孟斐斯妇人向她的丈夫指点我，要他买下我；因为她说：\Quote{我听说他们在卖他。}她便派一个太监到依市玛耳人那里，要求他们卖我；但那个太监不愿与他们交易，就回去了。那太监试探他们之后，便向女主人禀报说，他们为他们的奴仆索价甚高。她又派了另一个太监，说：\Quote{即使他们要两弥纳黄金，也要当心不要吝惜黄金；只管买下这少年，带他来这里。}他竟用了八十块金币买了我，却告诉女主人说用了一百块金币买我。我看见这事，便保持缄默，免得那太监受罚。\par

17. 你们看，我的孩子们，我忍受了多少大事，为的是不使我的弟兄们蒙羞。你们也要彼此相爱，以忍耐掩盖彼此的过错。因为天主喜爱弟兄的合一，以及一颗蒙悦纳的心所怀的爱。当我的弟兄们来到埃及，得知我归还了他们的银钱，并没有责备他们，反而安慰他们；在雅各伯死后，我更加爱他们，凡他所吩咐的，我都加倍遵行，他们就惊异了。因为我连最小的事也不让他们受苦；凡我手中所有的，我都给了他们。他们的孩子就是我的孩子，我的孩子就像他们的仆人；他们的生命就是我的生命，他们的一切痛苦就是我的痛苦，他们的一切疾病就是我的软弱。我的土地就是他们的土地，我的筹谋就是他们的筹谋；我并没有因世间的荣耀而在他们中间高傲自大，反而在他们中间如同最微小的一位。\par

18. 因此，我的孩子们，如果你们也遵行主的诫命，祂必在那里高举你们，并以美物祝福你们，直到永远。如果有人图谋加害你们，你们要以善行替他祈祷，你们就必蒙主救赎脱离一切凶恶。因为，看，你们见到我因忍耐竟娶了我主人的女儿。与她同来，还给了我一百塔冷通金子；因为主使这一切事为我效力。祂也赐我美丽，如同花朵，超越以色列的美丽之人；祂保守我直到老年，仍强壮而俊美，因为我在各方面都相似雅各伯。\par

19. 我的孩子们，你们也要听我所见的异象。有十二只鹿在吃草，其中九只被分开且散落在那地，另外三只也是如此。我看见从犹大生出一位童女，身穿细麻衣，从她那里出来一只无玷的羔羊，在其左边仿佛有一只狮子；所有的野兽都冲向祂，但羔羊战胜了它们，毁灭了它们，并将它们践踏在脚下。因着祂，天使、人类和全地都欢欣。这些事将在它们的时期，在末日发生。因此，我的孩子们，你们要遵守主的诫命，尊敬犹大和肋未；因为从他们中将为你们兴起天主的羔羊，以恩宠拯救所有外邦人和以色列。因为祂的王国是永恒的王国，永不动摇；但我在你们中间的王国将如守望者的吊床，在夏季之后便不复存在。\par

20. 我知道在我死后，埃及人将磨难你们；但天主必为你们辩护，并将你们领入祂应许给你们祖先的那地。但你们要将我的骸骨与你们一同带上去；因为当我的骸骨被带走时，主必在光明中与你们同在，而贝里雅尔必在黑暗中与埃及人同在。你们也要把你们的母亲齐耳帕带上去，将她安放在彼耳哈附近，靠近赛马场，在辣黑耳旁边。他说完这些话，便伸开双脚，长眠了。整个以色列和全埃及都哀悼他，极其悲痛。因为他对待埃及人也如同自己的肢体，向他们施恩，在一切工作、谋划和事务中都帮助他们。\par

% structure: p0146 source=h2 target=section reason=relative-heading-rank; base=h2

\section{十二．—— 本雅明遗训：论纯洁的心}

1. 本雅明对自己儿子们的讲话记录，他活了一百二十岁之后，向他们宣布。他亲吻了他们，并说：正如依撒格在亚巴郎百岁时所生，我也是在雅各伯百岁时所生。由于辣黑耳在生我时去世，我没有母乳；因此，她的侍女彼耳哈哺育了我。因为辣黑耳在生了若瑟之后，又过了十二年不生育；她禁食十二天祈求主，便怀孕生了我。因为我们的父亲深爱辣黑耳，祈求能见到两个儿子从她而生；因此，我被称作\Quote{日子的儿子}，即本雅明。\par

2. 因此，当我进入埃及时，我的哥哥若瑟认出了我，他对我说：\Quote{他们在我父亲面前说了什么，以致把我卖了？}我对他说：\Quote{他们把你的外衣蘸了血，送回去说：\InnerQuote{看看这是不是你儿子的外衣。}}他接着对我说：\Quote{正是这样，兄弟。因为当依市玛耳人带走我时，其中一人剥去了我的外衣，给了我一条腰带，鞭打我，并命我奔跑。当他去藏匿我的衣服时，一只狮子遇见了他，把他杀了；于是他的同伴害怕了，把我卖给了他们的同伴。}\par

3. 因此，我的孩子们，你们要爱天上的主天主，遵守祂的诫命，效法那善良圣洁的人若瑟；并要使你们的心意向善，正如你们所认识的我。凡是心意善良的人，就能正确看待一切。你们要敬畏主，爱人如己；即使贝里雅尔的邪灵引诱你们陷入各种烦扰的邪恶中，也不可让任何烦扰的邪恶辖制你们，正如它不能辖制我的哥哥若瑟一样。有多少人想要杀他，天主却保护了他！因为敬畏天主并爱人如己的人，不能被贝里雅尔空中之灵所击打，因为他被对天主的敬畏所保护；也不能被人或野兽的计谋所支配，因为他因他对邻人的爱而蒙主扶助。因为他甚至恳求我们的父亲雅各伯为我们弟兄祈祷，求主不将他们针对若瑟所图谋的恶事归到他们身上。于是雅各伯喊道：\Quote{我的孩子若瑟，你战胜了你父亲雅各伯的慈怀。}他拥抱他，亲吻了他两个小时，说：\Quote{在你身上将应验那关于天主的羔羊、即世界救主的预言，那无玷的羔羊将为违法者而被交付，那无罪的将为不敬虔的人而被处死于盟约的血中，为外邦人和以色列的救恩，并要消灭贝里雅尔和侍奉他的人。}\par

4. 我的孩子们，你们要知道善人的终局吗？要怀着善心效法他的怜悯，好使你们也能戴上荣耀的冠冕。善人没有邪恶的眼目，因为他向众人——即使他们是罪人，即使他们对他图谋恶事——都施以怜悯。因此，行善的人以善遮蔽了恶，并蒙那善者所护佑；他爱义人如同自己的灵魂。若有人受光荣，他不嫉妒；若有人致富，他不妒羡；若有人英勇，他赞美他；他信赖并颂扬那持守清醒的人；他向穷人施怜悯；他对弱者怀善意；他歌颂天主；他保护敬畏天主的人如同盾牌；他帮助那爱天主的人；他劝诫并挽回那拒绝至高者的人；他爱那拥有善神恩宠的人，如同自己的灵魂。\par

5. 我的孩子们，如果你们有善心，那么恶人也会与你们和睦相处，放荡的人会尊敬你们并转向善；贪婪的人不仅会停止他们过度的欲望，甚至会将他们贪婪的果实给予受苦的人。如果你们行善，就连不洁的神也会逃离你们；是的，连野兽也会因惧怕而逃离你们。因为当对善工的敬畏临到心智时，黑暗便从他那里逃遁。如果有人伤害一位圣人，他会悔改；因为圣人对他的辱骂者表示怜悯，并保持沉默。如果有人出卖一个义人的灵魂，而这义人虽祈祷，却暂时谦卑，但不久之后他会显得更加荣耀，正如我的兄弟若瑟一样。\par

6. 善人的心智不在\OriginalTerm{贝利亚尔}{Beliar}之灵的欺骗权势之下，因为和平天使引导他的灵魂。他不贪婪地注视可朽坏之物，也不为享乐的欲望积蓄财富；他不以快乐为乐，不伤害邻人，不纵情饮食，不因眼目的骄傲而犯错，因为主是他的份。善心不接受人的荣耀或羞辱，也不知晓任何诡诈或谎言、争斗或辱骂；因为主居住在他内，光照他的灵魂，他时时向众人欢喜。善心没有两样的舌头，即祝福与咒骂、侮辱与尊崇、忧愁与喜乐、安静与烦扰、虚伪与真理、贫穷与富足；而是对众人有纯全无玷的心意。它没有双重的眼光，也没有双重的听觉；因为他所做、所说、所见的每件事，都知道主鉴察他的灵魂，他洁净自己的心智，免得被天主和人定罪。但\OriginalTerm{贝利亚尔}{Beliar}的每一工作都是双重的，没有单一性。\par

7. 所以，我的孩子们，你们要逃避\OriginalTerm{贝利亚尔}{Beliar}的恶行；因为它给服从它的人一把剑，而剑是七种邪恶之母。首先，心智通过贝利亚尔怀胎，先是嫉妒；第二是绝望；第三是困苦；第四是掳掠；第五是贫乏；第六是烦扰；第七是荒凉。因此，加音也被天主交付于七种刑罚，因为每百年主就降一灾给他。他受苦二百年，在第九百年时在洪水中被毁灭，是为了他义人兄弟亚伯尔的缘故。加音受审判七百年，\OriginalTerm{拉默客}{Lamech}受审判七十个七倍；因为那些如同加音一样嫉妒与仇恨弟兄的人，将永远受同样的刑罚。\par

8. 所以，我的孩子们，你们也要逃避恶行、嫉妒和仇恨弟兄，而持守良善与爱。那在爱中有纯洁心智的人，不会为了淫乱而注视女人；因为他心中没有玷污，因为圣神安息在他内。正如太阳照耀在粪土和污泥上却不被玷污，反而使它们干燥并驱散臭气；同样，纯洁的心智被地上的污秽所包围时，反而建造，自己不受玷污。\par

9. 现在，我从义人\OriginalTerm{哈诺客}{Enoch}的话推测，在你们中间也会有恶行：因为你们将犯索多玛的淫乱，除少数外都要灭亡，并在女人中倍增过度的情欲；主的国必不在你们中间，因为祂要立刻夺去它。然而，天主的圣殿要在你们的土地上建造，并在你们中间得荣耀。因为祂要取它，十二支派和所有外邦人都要聚集在那里，直到至高者藉祂独生子的眷顾赐下祂的救恩。祂要进入圣殿的前面，主要在那里受凌辱，并被高举在木头上。圣殿的幔子要裂开，圣神要如倾泻的火降在外邦人身上。祂要从坟墓中复活，并从地上升到天上；我知道祂在地上将如何卑微，在天上如何荣耀。\par

10. 当\Person{若瑟}在\Place{埃及}时，我渴望看见他的面容和形貌；藉我父亲\Person{雅各伯}的祈祷，我在白昼醒着时，见到了他完整而完美的形象。所以，我的孩子们，要知道我快要死了。因此，你们各人要同邻人履行真理和正义，并秉公施行忠信的行为，遵守主的法律和祂的诫命；因为这些是我教给你们取代一切遗产的。你们也要将这些交给你们的孩子们作为永久的产业；因为\Person{亚巴辣罕}、\Person{依撒格}和\Person{雅各伯}都是这样做的。他们将这些都赐给我们作为产业，说：要遵守天主的诫命，直到主向万民显示祂的救恩。那时，你们将看到\Person{哈诺客}、\Person{诺厄}、\Person{闪}、\Person{亚巴辣罕}、\Person{依撒格}和\Person{雅各伯}，喜乐地从右边复活。那时，我们也要复活，每人管理我们的支派，朝拜天上的君王，祂曾以谦卑之人的形像出现在地上。凡在地上相信祂的人，都要同祂一起喜乐；然后，所有人将复活，有的得光荣，有的受羞辱。主将首先审判\Place{以色列}，因为他们对祂行了恶；因为当祂作为拯救者出现，天主降生成人时，他们不信祂。然后，祂将审判所有外邦人，就是那些当祂出现在地上时不信祂的人。祂要在外邦人的选民中责备\Place{以色列}，正如祂在\Place{米德杨}人中责备\Person{厄撒乌}一样，\Person{厄撒乌}曾欺骗他的弟兄，使他们陷入奸淫和偶像崇拜；他们就远离了天主，在敬畏主的人的分中，如同不是子女一样。但若你们在主面前行走于圣洁中，你们将再次在我内居于希望中，全\Place{以色列}将被聚集到主那里。\par

11. 我不再因你们的劫掠而被称为掠夺的狼，而是主的工人，将食物分给行善的人。在末世，将从我的后裔中兴起一位蒙主宠爱者，在地上听到祂的声音，以新知识光照所有外邦人，带着知识的光冲入\Place{以色列}以取得救恩，像狼一样从它那里夺走，并赐予外邦人的会堂。直到时代的终结，他将在外邦人的会堂中，在他们的统治者中，如同众人嘴中的乐曲；他的工作和他的话语将被记录在圣书中，他将永远是天主所拣选的；因着他，我父亲\Person{雅各伯}教导我说：他将填补你支派所缺乏的。\par

12. 当他讲完这些话后，说：我嘱咐你们，我的孩子们，将我的骨骸从\Place{埃及}带上去，葬我在\Place{赫贝龙}，靠近我的祖先。\Person{本雅明}寿享一百二十五岁，善终，他们将他放入棺材。在\Place{以色列}子民出\Place{埃及}后的第九十一年，他们和他们的兄弟秘密地将他们祖先的骨骸带到一个名叫\Place{客纳罕}的地方；并将他们葬在\Place{赫贝龙}，靠近他们祖先的脚旁。他们从\Place{客纳罕}地返回，住在\Place{埃及}，直到他们出\Place{埃及}的那一天。\par

\SourceRecord{Translated by Robert Sinker.；Ante-Nicene Fathers；Vol. 8.；Edited by Alexander Roberts, James Donaldson, and A. Cleveland Coxe.；Buffalo, NY: Christian Literature Publishing Co.,；1886.；<http://www.newadvent.org/fathers/0801.htm>.}
